\chapter{燃烧的革命}

\section{中央苏区的党}


中央苏区的发展和军事推动直接相关，武力是中央苏区创建、发展的第一推动器，这不是什么革命史话语的表述，而是当年历史的现实。因此，无论是中共的组织建设还是苏维埃政权的创立，都与军事的占领密不可分，正如红军大举进入赣西南地区后中共报告所谈到的：“党的组织的发展，是红军打来之后才发展的，毛泽东起草计划，要三天发展十万党员。”\footnote{《赣西南会议记录（1930年10月13日）》，《中央革命根据地史料选编》（中），第628页。}虽然三天发展十万党员的说法有些夸张，却是政治力量通过军事占领大力推动的真实写照。


中共是强调组织严密的革命党，党的组织在革命中发挥着核心作用，因此，建构强大的党组织是中共强化自身的必要举措。政党的组织力量虽然不能单纯用数量衡量，但一定程度的人员覆盖仍然是必要条件之一，就此而言，中共在控制区域内的数量发展远非国民党所可比拟。作为全国的执政党，国民党在南京国民政府成立后并没有迅速发展，相反在两广地区由于清党的影响还出现萎缩。就江西而言，1928、1932、1934年国民党员人数分别为10869人、16513人、15410人，增长速度相当缓慢。1935年与人口比仅为1∶722。\footnote{转见王奇生《清党以后国民党的组织蜕变》，《近代史研究》2003年第5期。}赣南的会昌县当时拥有20多万人口，1929、1930、1931年国民党员人数却分别只有216人、189人、140人，1932年后因会昌苏维埃化，国民党员人数更是急剧减少。\footnote{《民国时期会昌的国民党等党团组织概况》，《会昌文史资料》第3辑，政协会昌文史资料编辑委员会编印，1989，第53页。}1933年底，江西邻近省会南昌的丰城、清江两个大县分别只有国民党员250人、210人，党务工作“无进展，甚至已陷停顿状态”。\footnote{《视察丰城清江新淦三县报告书》，《军政旬刊》第5期，1933年11月30日。}根据同年省会南昌的调查，国民党员也由1920年代末的两千多人减少到六百来人，且“没有方法能把这几百个党员团结组织起来”。\footnote{《中国国民党南昌市党部改组三月来的工作》，南昌《市政半月刊》第1卷第5、6期合刊，1934年10月1日。}不仅如此，党员的来源和结构也不合理，代表性大有问题。江西1934年的国民党党员分类统计中，农业口（含林渔牧）党员共计900人，其中业主668人；工业口党员共计56人，其中厂主27人，职员29人。\footnote{《中国国民党江西省党务统计报告》，中国国民党江西省执行委员会编印，1934，第21页。关于南京国民政府建立后国民党组织具体状况，可参阅王奇生《党员、党权与党争》（上海书店出版社，2003）有关论述。}有产阶层在国民党员中占据绝对优势，占人口绝大多数的普通农民和工人比例甚低。


国民党当时是奉行“以党训政”的非竞争性政党，即便如其所说，以党训政不是党员训政，而是以党义训政，党员仍然应该是其发挥统治力的一个不可或缺的环节。在党员的数量和代表性都无法在社会上形成有效影响时，国民党组织功能的发挥自是困难重重。所以蒋介石在江西曾批评道：“现在许多党部和党员……领了经费，还不能努力剿匪工作，以增加我们剿匪的效能。那么只成就了一个纯粹的党的衙门，完全是一个假门面，失却了我们党员和党的资格。”\footnote{蒋介石：《剿匪区县长及党政人员的职责与行政的方法》，《先总统蒋公思想言论总集》第11卷，第246页。}这道出了南京政府统治下国民党党组织的实况。


和国民党相比，中共对党的建设可谓全力以赴。其经典的表述是：“巩固党本身的组织，坚强党的无产阶级基础，最高限度的提高各级党部——从支部起——的积极性，严紧党的纪律，增加党的领导作用，成为最先决最重要的问题。”\footnote{《党的建设问题决议案（1931年11月）》，《中央革命根据地史料选编》（上），第635～636页。}对组织的高度重视使中共党员无论是数量和质量都比国民党明显高出一筹。即便是在中共尚未控制大片根据地的1928年，江西全省已经拥有中共党员4000余人，包括14个县委、1个市委、46个区委、118个支部。\footnote{《江西工作近况（1928年7月3日）》，《江西革命历史文件汇集（1927～1928年）》，第261页。}中央苏区逐渐发展起来后，中共组织在江西急剧发展。1932年3月，苏区归属江西省辖的16个县全部建立了县委，成立124个区委、998个支部，党员发展到2.3万余人。会昌作为新占领地区，中共党员人数也达到660人，数倍于此前国民党员的数量。组织发展较为健全的公略县党员与人口比已经超过1∶30。\footnote{《江西苏区中共省委工作总结报告》，《中央革命根据地史料选编》（上），第490页。}1932年7月，仅对公略、兴国、胜利、于都、宁都、乐安、永丰7县的统计，党员总数就达到13124人，其中兴国一县达4063人。\footnote{《中央苏区组织统计表（1932年7月）》，《中央革命根据地史料选编》（上），第661页。}这已经接近国民党在江西全省的党员数。此后，中央苏区又开展大规模发展党员的冲锋运动，1933年6月，全省党员67904人，9月，党员达到97451人。\footnote{《党的组织状况》，《中央革命根据地史料选编》（上），第674页。}这实际还只是中央苏区划定的江西省范围，不包括闽浙赣、湘鄂赣、湘赣苏区。


福建中共党组织集中在闽西。1930年3月的统计显示，福建省有党员1318人，其中永定、龙岩各有300人，上杭100人，长汀、武平分别为15人、35人，这几县的党员人数占到全省中共党员数的一半以上。\footnote{《福建组织状况统计表（1930年3月8日）》，《福建革命历史文件汇集》甲4册，第146～153页。}随着闽西苏区的开辟、发展，闽西中共党组织迅速壮大，1932年3月党员人数达到6800人左右，1933年6月达20000人。闽西、赣南加上军队中几近半数的党员，整个中央苏区党员最多时应不低于15万人。\footnote{1934年8月，中共驻共产国际代表团特别致函共产国际执行委员会，就中共党员人数问题作出解释，提到中共当时在国内拥有40余万名党员，其中江西、福建两省分别有党员10万、4万人。见《中共驻共产国际执行委员会代表团给共产国际执行委员会政治书记处政治委员会的声明》，《共产国际、联共（布）与中国革命档案资料丛书》第14册，第203～204页。}以300万总人口计，党员在人群中的比例占到5%。


一定数量的党员是组织发挥作用的必要前提，党员质量则是决定组织坚强与否的要素。中共党的组织原则直接承袭自苏俄，纪律严明、令行禁止，强调牺牲和奉献精神。这一原则在中央苏区的中共建设中，基本得到贯彻。中共对发展党员态度认真，红四军第九次代表大会决议中提出的入党条件是：“1.政治观念没有错误的（包括阶级觉悟）。2.忠实。3.有牺牲精神，能积极工作。4.没有发洋财的观念。5.不吃鸦片，不赌博。”\footnote{《中国共产党红军第四军第九次代表大会决议案》，《中共中央文件选集》第5册，中共中央党校出版社，1990，第813页。}1929年7月，红四军第三纵队政治部编印的《党员训练大纲》中，对发展新党员的步骤作出详尽提示，要求介绍人在选择介绍对象时，要从如下几方面加以考察：“从他的家庭经济背景，考察是否有革命之须要”；“从他平时做事待人，看他是否忠实可靠”；“从斗争中看他是否勇敢不怕得罪人”；“从他的交友中或反对者，各方面看他是否好人”；“在他谈话中看他能否守秘密”；“从他的脾气上看他是否会服从”；“从他平常看书上看他的思想是否革命”。\footnote{《党员训练大纲》，红四军第三纵队政治部编印，1929。}经过严格考察之后，若基本符合入党条件，即由党组织循谈话、讨论、通过等程序予以发展。虽然在苏区农业社会及军事紧张的现实环境下，中共的无产阶级理念未必能够顺利贯彻，组织条件也难以完全满足，但强调忠实可靠、守秘密、服从等基本品质，遵循可操作的发展程序，还是有助于保证党的凝聚力和纯正性。


和国民党基层组织多为放任自流的状况比，中共十分重视基层组织尤其是支部的建设，1932年3月底，兴国等15县共成立支部998个，差不多每乡有一个支部。\footnote{《区委支部统计表（1932年3月底止）》，《中央革命根据地史料选编》（上），第661页。}中共各级机关对加强支部建设的阐述不胜枚举。从1931年中共赣西南特委颁布的支部工作条例中，可以看出中共对此所作计划的具体和详尽。文件规定：“有三个同志以上即可成立支部”，“支部在十人以上就要分小组”，支部成立支部委员会。支部委员会每七天开会一次，小组会每五天开会一次，支部党员大会每半个月召开一次。支部会议“多讨论支部本身工作和实际工作以及工作方法，每次会议均须讨论中心问题”。有意思的是，文件还提到：“会场秩序不必太庄严亦不必太散漫，同志有时可以嬉笑，但要有限制，以不妨碍工作大家的注意力为原则。”\footnote{《中共赣西南特区委关于支部工作的部署（1931年6月7日）》，《东固·赣西南革命根据地史料选编》第2册，中央文献出版社，2007，第709～717页。}对会场气氛这样看似细枝末节却相当程度上可以影响会议效果的问题也不放过，充分显示出中共严谨细致的一贯作风。


对党员的思想教育是中共保持和发挥自身特质的必由之举。随着苏区的发展，各级党员教育机构次第设立。1932年春，中共苏区中央局开办中央党校。1933年春起，中央苏区开办了一系列正规干部学校，主要有马克思共产主义大学（即中央党校）、苏维埃大学、列宁团校等。同时，开办各种短期训练班，随时随地开展对党员的教育、训练。江西永新1934年6月设班培训的干部（包括任支部书记、县区常委及开展新区工作等）达104人。\footnote{赵可师：《赣西收复区各县考察记》（四），《江西教育旬刊》第10卷第8期，1934年8月11日。}1933年8月，中共中央组织局专门就党内教育计划问题致信各级党部，要求从中央到省委、县委，将举办各类短期圳练班作为培训干部的一项重要措施。省委训练班负责训练造就县一级干部、巡视员、县委训练班的教员及区一级主要干部，训练时间以四星期为限。县委举行流动训练班训练区一级干部和部分支部书记及支部流动训练班的教员。这些经过培训的干部充实到各级机关后，对提高干部队伍的素质及其忠诚度，严密组织，加强中共党组织和干部体系的领导能力，有着十分重要的作用。经过当年干部学校训练的红军干部回忆：

\begin{quote}

当时的政治课内容既讲马列主义又讲苏联的十月革命，也有一些政治工作常识，既有如何在连队中做好工作，也有如何做好党的群众工作方面的知识和课程。这次学校虽然不长，但对我的一生产生了十分重要的影响。不仅使我真正学到或懂得了一些无产阶级革命的目的和意义，而且也是我人生的重要转折点。\footnote{《罗元发回忆录》，光明日报出版社，1995，第56页。}
\end{quote}


当然，像中央苏区党这样所谓在深山中成长起来的政治力量，先天不足、后天难调的缺陷也客观存在。苏区教育文化落后，民众很难从理论上了解中共的思想体系，为保证党的力量发挥，又必须要求短期内达到一定的党员覆盖，因此对党员的质量很难提出更高的要求，出身贫寒、忠诚老实几乎是可以悬揭的高标准了。共产国际和中共中央虽然一再强调要把苏区党锻造成为一个无产阶级政党，但在当时这未免有点不切实际。尤其1932年后，苏区超常发展党员，对中共组织严密度的冲击更大。1933年5月，苏区开展红五月征收党员运动，一个月内征收了两万多名党员。这些党员基本都是由各地列出指标，然后采取速成办法发展：“十五个县的报告所写明的，只有博生、赣县、瑞金部分的区和寻乌的澄江办了新党员训练班，瑞金部分的用开会的方式训练了一天。这就看到对新党员入党时给以基本训练的工作，是极少注意到的。”\footnote{罗迈：《红五月征收党员的结果与教训》，《斗争》第20期，1933年8月5日。}更极端的还有：“博生田头区山头王支部质问群众‘你为什么不加入党，难道你是反动派么？’以至如石城坝口区某乡支部派了两个同志拿一本簿子一支铅笔在各村各屋填名字，口里喊着‘加入共产党’。”这样发展来的党员，对党的认识模糊，江西宜黄“大多数党员不明了党的组织，往往党的组织与群众组织分不清楚，如宜黄吴村区出席县党组织代表大会的代表是苏维埃介绍来的，并介绍非党员来充当县党代表大会做代表”。\footnote{《党的组织状况》，《中央革命根据地史料选编》（上），第678、697页。}同时，组织纪律也无法保证，江西会昌反映：“有支部连开会都召集不到来开。”\footnote{《中共会昌县委十、十一两月工作报告》，《江西革命历史文件汇集（1932年）》（二），第373页。}更荒唐的是，江西南广“有些党员不承认加入了党（头陂、城市、白沙、巴口、长桥，尤其是长桥区有半数不承认加入党）……一般的是不了解为什么加入党，党是什么（巴口有党员说加入党不打路条不作挑夫）”。\footnote{《中共南广县委给苏区中央局的综合报告（1932年10月）》，《江西革命历史文件汇集（1932年）》（二），第92页。}因此，一些党员遇到风吹草动，就可能出现问题，1933年初，国民党军进攻寻乌，其间全县653名党员中，反水的达72人。\footnote{《中共寻乌县委一个半月动员工作总结报告》，《江西革命历史文件汇集（1933～1934年）》，第36页。}


短时间内在狭小地区把党员数量发展到十多万人，上述问题的出现几乎是不可避免，这是中共急于在苏区建构庞大组织网络必然付出的代价，有利则必有其弊。关键是，尽管有不少问题，大发展的结果，还是使中共短时期内最大限度地发挥出了自身的影响力，而中共空前的组织力和严密度，虽有被稀释的现象，尚不足以遭到根本动摇。这主要是因为除党员外，中共更培养了一支高质量的干部队伍，事实上，中共当年所体现出的强大组织力，更多的还是依靠中低级干部这样一个可以有效上传下达的支撑力量得以实现。


干部是中共组织的坚强骨骼。与中共创建、发展过程中知识阶层所具领航者地位一样，赣南、闽西早期党组织中，地方上一般都是接受马克思主义思想的知识分子起着核心作用。1928年12月，江西省委正式委员13人，其中知识分子占到8人，候补委员7人中知识分子占5人。\footnote{《江西省组织状况一览表（1928年12月）》，《江西革命历史文件汇集（1927～1928年）》，第316页。}1929年8月，闽西苏区各县县委常委统计成分的共23人，其中知识分子14人，农民7人，工人2人，知识分子占60%。\footnote{《各县县委调查表》，《福建革命历史文件汇集》甲8册，第113页。}何友良对早期东固苏区的研究显示，东固领导群体的20人中，青年学生、知识分子多达17人。\footnote{何友良：《革命源起：农村革命中的早期领导群体》，《江西社会科学》2007年第3期。}这些知识分子干部投身革命多出于自身的理性选择，对主义、革命怀抱信仰，有强烈的献身精神，为中共实现对地方的控制起到不可替代的作用。1932年任泰宁溪口区区委书记的钟国楚回忆，他当时开展工作较好的重要原因是：“有当地干部，陈家源的支部书记陈国夫，从当时来看他是个知识分子，工作有办法，还有蒋坊一个知识分子，是团支部书记还是党支部书记，一时记不清。这两个同志……是我们工作中的左右手。”\footnote{《南京军区顾问钟国楚同志口述记录稿》，《泰宁文史资料》第1～3辑合订本，第14～15页。}


随着苏区的发展，苏区党员的源流日渐丰富，干部成分也相应发生变化。1929年6月，江西各县县委委员中，知识分子占到60%。\footnote{《江西省组织状况统计表（1929年6月6日）》，《江西革命历史文件汇集（1929年）》（一），第235页。}1931年10月湘赣苏区统计，全苏区县委一级的干部共146人，其源流为：产业工人3人，手工业工人28人，苦力工人3人，店员工人3人，雇农10人，贫农64人，中农10人，兵士2人，知识分子20人，其他3人。\footnote{《中共湘赣苏区省委综合工作报告》，《湘赣革命根据地》（上），第110页。}作为强调无产阶级成分的结果，工人、雇农的比例明显上升，知识分子比例则下降到13.7%。应该说，这两种成分的比例都超过了其在苏区中实际人口的比例，前者代表的是中共的阶级立场，后者则反映的是知识分子在发起社会革命时难以替代的地位。


作为一个工人阶级政党，大规模发展普通工农尤其是工人入党，提拔他们充实进领导部门，是中共明确自己阶级属性的必然之举。为此，中共投入相当精力，开办党校、举办训练班，提高干部素质。从1932年中共河西道委一则开办短训班的通知，大致可看出中共在这方面的努力：

\begin{quote}

目前工作日益开展，新区域工作极需干部工作，因此，道委决定开办短期训练班，时间两礼拜。望即根据下列指示送学生来，至要！

1. 送来的学生，必须是现在区一级工作或担任支部书记的同志，成份要是工人、雇农、贫农，而且要是本地同志，最好稍能识字的。

2. 送来的学生必须身体强健，有活动能力和学习精神，能听党调动工作，并无反革命派别嫌疑的。\footnote{《中共河西道委通知（1932年4月20日）》，《江西革命历史文件汇集（1932年）》（一），第90页。}
\end{quote}


经过中共的努力，一批新干部迅速成长。1933年5月的统计，江西县一级干部中，工人成分占到46%，贫农成分占到44%，其他所有成分只占10%。\footnote{罗迈：《把提拔新的干部当作组织上的战斗任务》，《斗争》第25期，1933年9月5日。}1934年1月第二次全国苏维埃代表大会835名代表中，源自产业工人者8人，手艺工人244人，苦力工人53人，店员工人12人，雇农122人，贫农303人，中农25人，商人4人，其他64人。\footnote{《红色中华》第145期，1934年2月3日。}工人、雇农、贫苦农民成为干部的主力。1933年夏中共中央局对江西省16个县县一级干部所作调查显示，419名干部中，1927年以前入党的13人，占3%；1928～1929年入党的52人，占12%强；1930年入党的125人，占30%；1931～1932年入党的190人，占45%；1933年入党的39人，占10%弱。参加过游击战争或暴动的185人，占44%；到县工作之前在支部和区一级工作过的占81%。\footnote{《党的组织状况》，《中央革命根据地史料选编》（上），第687页。}可以看出，苏区干部形成了一个梯级结构，以参加过暴动或游击战争的干部为核心，同时大批提拔新生代干部，实现新老结合、新陈代谢。


中共很注意对干部的宣传教育，苏区时期中共高中级干部多为怀抱着理想投身革命，有很强的献身精神，工作作风和思想情操都令人印象深刻。黄克诚回忆，当年他由于不赞成攻打中心城市而被批评为右倾，与同事间发生不少争论：“我同军政治委员贺昌一起走，他继续批判我的右倾机会主义，我不服，就跟他争论。部队宿营时，我俩还是住到一块，继续争论，吵得很厉害，彼此各持己见，谁也说服不了谁。我对贺昌表示，准备同他争论二十年。贺昌不愧是个真正的共产党人，他作为上级，我无论怎样同他争吵，他都不在乎。争吵归争吵，吵过之后，照样相处，毫不计较，也不影响工作。”\footnote{《黄克诚自述》，第108页。}一大批能将个人情感、私利和工作、事业清楚分开的干部的存在，是中共组织和领导力量得以顺利发挥的最重要保证。


作为一个还在为自己生存权利奋斗的政党，苏区时期的中共干部很少沾染官僚作风，工作认真、能吃苦、深入实际是他们的共同特点。担任中央妇女部长的周月林回忆：“我们作报告一般都是先列个提纲，到下面去了解一些情况，就可以结合着讲了。讲完后对下面工作存在着什么障碍，应该怎么办，提出来大家讨论，拿出解决的办法。”\footnote{《女英自述》，江西人民出版社，1984，第155页。}曾任中共长汀县委书记的李坚真谈到她当年的工作体会：

\begin{quote}

为了管好一个县，当好这个“家”，我根据自己的特点，依据的是“三勤”，也就是腿勤、手勤、嘴勤。“腿勤”就是经常往下面跑，接触实际，联系群众，发现问题及时处理。无论在省委或县委工作，我都很少住在机关，靠的是一双赤脚板走遍各区、乡、村，和群众住在一起，吃在一起。

……

“手勤”就是走到哪里，我就在哪里拿起锄头，挑起扁担和群众一起劳动。下田插秧、割禾，上山砍柴，挑水煮饭，喂猪样样都干。群众把我当成自己人。

……

“嘴勤”就是多宣传……我从小特别爱唱山歌，闽西群众也和我们广东客家人一样喜欢唱山歌。我就结合形势、任务编些山歌和群众一起唱，通过唱山歌提高群众的阶级觉悟。\footnote{《女英自述》，第163～164页。}
\end{quote}


千百年来习惯了高高在上的统治者的普通民众，骤然见到这样的“县官”，其反应可想而知，而高中级的干部以身作则，又会对基层干部形成潜移默化的影响。中央苏区干部的良好口碑就是由此酿成的。正是有这样一批富有献身精神的干部的存在，使中共即使在组织快速膨胀的形势下，也能保持一个坚强的核心和集中的凝聚力。


虽然中共提拔了一大批工农干部，但苏区时期干部的核心仍然是经过数年革命洗礼、具有坚强信念和革命精神的知识干部，知识者所具的这种地位，在世界革命史上都是通例，共产革命也不例外。知识者得风气之先，革命尤其是共产革命又充满着理想和浪漫，共产主义的理念对当年知识阶层有着无与伦比的吸引力，成为凝聚这一批革命核心的坚强精神力量。不过，共产革命是无产阶级革命，保证无产阶级对革命的领导权在共产革命的话语中意义非同寻常。而知识阶层作为一个整体充其量只能作为革命的同路人，因此无论是苏俄和共产国际，还是本身就出身于知识阶层的中共领导人，对于知识阶层在革命中的地位都有一种缘于阶级分析的担忧。体现着中共领导人的信仰和热情，他们对自身曾经隶属的那个阶层多采取无情的批判和检视态度，虽然多年后的经验告诉我们，个体乃至阶层的理念受到多种因素的影响，并不一定和其社会政治地位发生必然的联系，但当年的逻辑自有其不可移易的权威。所以，仔细考察中共革命的轨迹，可以发现不同时期不同地域知识阶层的命运大体类似，区别只在于政策执行中的刚性和柔性而已。


1932年中共中央大批干部到达中央苏区后，对知识阶层执行偏于刚性的政策，具体而言，就是以排除知识阶层干部为实际结局的“唯成分论”盛行：

\begin{quote}

在干部路线上大搞唯成分论，过分强调红军领导骨干必须是无产阶级成分，无产阶级分子，向各地区各部队派遣大批“钦差大臣”，去进行所谓“改造和充实各级领导机关”，搞乱了干部队伍。当时部队绝大多数同志是农民出身，而他们却只提拔工人出身的人。不问其是否具备干部条件，只要是无产阶级成分的就提拔。\footnote{《李志民回忆录》，第221～222页。}
\end{quote}


这样的批评得到了许多当事者回忆的印证，应该不是空穴来风。就连毛泽东当时也忍不住抱怨，一些人对待干部：“普遍的只讲成分，不讲工作，只要是出身坏，不管他有怎样长久的斗争历史，过去与现在怎样正确执行党与苏维埃的路线政策，一律叫做阶级异己分子，开除出去了事。”\footnote{毛泽东：《查田运动的初步总结》，《斗争》第24期，1933年8月29日。}


毛泽东提到的这种现象，在中共中央要求保证党的无产阶级纯洁性口号下，确实屡见不鲜。在成分论影响下，苏区成分高的干部普遍不被信任，纷纷遭遇被洗刷的命运。胜利县1932年11月“开除了二十多个同志，开除的原因，大多数是富农分子”。\footnote{《中共胜利县委十一月份工作报告大纲（1932年12月10日）》，《江西革命历史文件汇集（1932年）》（二），第382页。}1932年11月至1933年8月，江西乐安洗刷县、区干部可确认成分者12人，其中地主、富农11人，贫农只有1人，有4人仅仅是因查田运动中被定为地主成分而遭洗刷，其他的罪名也多是所谓参加勇敢队、扯布告等。\footnote{《党的组织状况》，《中央革命根据地史料选编》（上），第696页。}闽浙赣省委报告，经过1932年初的肃反，“党的干部已大大的撤换了……全省的知识分子已去十分之九还要多一点”。\footnote{《中共闽浙赣省委向中央的报告》，《闽浙皖赣革命根据地》（上），第561页。}通过洗刷，知识分子出身干部数量已经微乎其微，而工人和贫农出身者占到干部队伍90%左右，成分和出身主导着干部的命运。


中央苏区处于教育文化比较落后的山区，要造就一支文化素质较高的干部队伍，在这里先天就存在较大困难。而在强调成分政策影响下，大批洗刷旧干部，提拔新干部后，这一不足表现更为明显。1933年年中的调查，江西县一级干部中，1931～1933年入党的占到总数的近55%，其中于都、瑞金两县“只满一年或不满一年党龄的干部在两县的领导机关中占比较的多数”。这些人中，能写东西的占37%弱，不能写的占63%，完全不识字的达到20%。\footnote{罗迈：《把提拔新的干部当作组织上的战斗任务》，《斗争》第25期，1933年9月5日。}县一级干部尚且如此，区、乡一级更可想见。福建长汀有些区委书记报告工作时，“涨红了脸，什么也报告不出”，问他们如何开展工作，“只答复‘我不识字’，‘决议又看不懂’就完了”。\footnote{荣光：《掩盖消极怠工的好办法》，《青年实话》第16号，1932年4月25日。}不少的笑话也由此产生：“在长汀训练班上，有一个支部书记，我们问他共产青年团是什么？答：‘红军的铁军’。”\footnote{《团闽粤赣省委三个月工作报告》，《闽粤赣革命历史文件汇集1932～1933》，第175～176页。}“许多问题不要说很少人到群众中去宣传，连宣传人员自己也不明白。据说李卜克内西及卢森堡纪念日，宣传员去宣传，群众问他‘李……卢……在哪里？’他说：‘在团部’。”\footnote{《丘泮林闽北巡视报告（1930年10月）》，《闽浙皖赣革命根据地》（上），第217页。}


文化水平的低下，不仅仅是文化问题，正如当时的领导者也意识到的：“当我们估计干部对党对革命的坚定性的时候，文化程度没有什么关系，但文化程度对于政治的发展和工作能力的进步则有很大的关系。”\footnote{罗迈：《把提拔新的干部当作组织上的战斗任务》，《斗争》第25期，1933年9月5日。}干部素质的下降比党员素质下降对中共组织更具负面意义。中央苏区后期，各级政权和组织出现一定程度的软化与此应不无关系。邓颖超记录了组织涣散的实际状况：瑞金城区党员大会原定下午1时开会，结果由于大家拖拉，一直到下午4时才开成会，出席人数仍只占应到会人数之半。会上，东郊支部书记反映：“支部决议，一般党员不执行，支部无检阅工作，一个多月未开一次支部会，支干会开了一次。发展党员，支部不知。”而南郊支部的书记则报告：“支部决议没有执行。赤卫军下一次操，少先队未下操。”\footnote{邓颖超：《瑞金城区党员大会经过与教训》，《党的建设》第4期，1932年9月10日。}部分基层组织实际处于半瘫痪状态。


其实，当中共中央强调所谓工人成分时，已不免郑人买履之嫌。中央苏区并无大工业，大部分工人成分其实是手工业者、乡间手艺人或雇农。1934年初，中央苏区拥有工会会员145000人，涵括了95%以上的工人，与整个苏区300万人口比，只占约4.8%。\footnote{《中华全国总工会给赤色职工国际报告（1934年3月1日）》，《中国工会历史文献》第3册，工人出版社，1958，第624页。}他们在苏区农村所占比例小，其生活、思想方式也和普通农民并无二致，不少人仅因为成分的关系，得以在党、政权、军队中占据重要地位。根据工会方面给职工国际的报告，1933年间，“在中央苏区，工会供给了差不多一万个工人干部到苏维埃、红军、党、团、各种群众团体工作。各苏区在苏维埃机关的负责人，工人占了三分之一到二分之一。各级苏维埃主席多数是工人、雇农、苦力，红军长官工人占了百分之五十以上”。\footnote{《中华全国总工会给赤色职工国际报告（1934年3月1日）》，《中国工会历史文献》第3册，第626页。}虽然为迎合共产国际对无产阶级政党性质的强调，报告可能夸大了工人成分的干部比例，但提拔工人干部的趋势应没有疑问，而在如此小的人口比例中要速成这么多的干部，结果可想而知。以致各地普遍反映：“工人斗争情绪不如农民（三期党校测验，百分之九十五的人是如此意见）。”\footnote{邓颖超：《实际为巩固与加强无产阶级领导权而斗争的检讨》，《斗争》第1期，1933年2月4日。}工人的特殊地位，甚至引起了普通农民与他们的冲突：

\begin{quote}

赣东北的横峰县乡下有个雇农，名叫陈克思，十六岁，由工会青工部的领导，同雇主订立了劳动合同。在劳动合同上规定，实行每日工作六小时，不担四十斤以上的担子，工钱从八元增到十六元。而这个“老板”却是贫农。另外有个十五岁的牧童名叫吴树德，每日工作四小时，工资由四元增到十元，挑担子不得超过三十斤，而“老板”又是贫农。\footnote{洛甫：《五一节与劳动法执行的检阅》，《斗争》第10期，1933年5月1日。}
\end{quote}


应该指出的是，中共中央对苏区阶级构成的现实也不是没有了解，对无产阶级成分的强调，更多是回应共产国际的担忧——在广大农民中成长起来的中国党的无产阶级性质。在问题日渐暴露后，中共中央开始作出反省和调整。1933年5月，苏区中央局针对前一段时间盲目排斥知识分子的现象作出决议，批评中央苏区在发展和巩固党的组织中“左”的错误，强调：

\begin{quote}

把那些完全准备着在共产党与共产国际的基础之上为无产阶级的目的而斗争的知识分子吸收到党里面来，对于无产阶级的事业，是有很大的意义的，尤其现在苏区的党，领导着工农民主专政的党，由于许多条件还不能及时的造就工人阶级的知识分子的时候，吸收革命的知识分子的入党，更是必须的。

在努力发展党员的任务上，应该把工人雇农苦力中所有的优秀分子和最革命的贫农吸收入党，因为这是党的基础，反对“共产党是穷人的党”的非阶级的口号，同时应该吸收真正进步的革命的知识分子和中农入党，纠正对于他们的关门主义。\footnote{《关于纠正发展和巩固党的组织中错误倾向的决议》，《中共中央文件选集》第9册，第202～205页。}
\end{quote}


同年12月，王明进一步指出：“因为中国苏维埃革命巨大的胜利，在小资产阶级的知识分子当中，发生了一种显然的变化。就是发生了倾向革命和共产党的左倾情绪。我们现在在这方面的任务，首先就是加紧在这些小资产阶级知识分子当中的工作，把他们之中许多人变成为反帝运动的积极战士，应用他们作为宣传鼓动员的力量，并且把他们之中的各种专门人材吸收到我们方面来去帮助中华苏维埃共和国的建设工作。”\footnote{王明：《革命、战争和武装干涉与中国共产党底任务》，《王明言论选辑》，第371页。}这是中共自1927年国共分裂后对小资产阶级政治倾向作出悲观判断后，重新对他们的政治倾向作出估定的重要起点，不仅是对知识分子阶层，实际包含着中共对中国社会阶级关系重新思考的内涵。此后，从政策层面上，知识分子被歧视、排斥现象得到了一定程度的遏制。


尽管尚有不如人意之处，但毋庸置疑的是，中共在中央苏区建立的党尤其是干部的组织体系，表现出了强大的控制力和生命力，中共在自己控制的地区内初试啼声，其成效已足骄人。

\section{中央苏区的政权}

\subsection{苏维埃政权的建立}


中共是一个有理想、有目标、有组织的革命政党，一开始就将建立新型政权作为革命的不二法门。国共合作领导的国民革命果实被蒋介石夺走后，面对国民党的武力镇压，中共别无选择地走上了武装革命之路。对于共产主义革命的倡导者而言，革命并不仅仅是控制权力的转移，而是社会政治的全新革命性转换，因此，建立新政权，实现社会结构的改造，是革命的题中应有之义，其悬揭的革命政权模式就是当时世界唯一的苏俄式苏维埃政权。共产国际代表描述了中国初期苏维埃政权的状况：“试图躲在自己的区域内，用万里长城将自己同外界隔开，建立一个摆脱赋税、摆脱地主统治等等的国家。”\footnote{《米特凯维奇的书面报告〈中国苏维埃经验〉》，《共产国际、联共（布）与中国革命档案资料丛书》第7册，第510页。}如果撇去其中的批评性表述，这一判断大致还是反映出中共苏维埃革命的追求。


传统中国社会政治组织松散，政府对地方控制力有限，绅权在地方社会影响重大，晚清绅权尤为伸张。民国时期这种状况仍基本得到延续，地方豪绅依然相当强势。根据抗战时期对78位县长的调查，他们最感困难的几件事中，财政困难居首位，其次即为土劣猖狂，有31人选择，占总数的39.74%，\footnote{企华：《县长认为最困难的几件事》，《新赣南旬刊》第4卷第8期，1942年12月10日。}可见地方土劣势力的强盛。中央苏区的赣南、闽西在各自的省份都处于较为偏僻山区，离行政中心省会较远，行政控制力不足，绅权尤为强盛。民国初年，广昌“一切县政，多有操在县绅之手”，县机关“大都有名无实，为私人分肥机关，绝少有为民众兴利除弊，实事求是者！甚有于民十五年，当革命军兴，适县长田瑞璜在任病故，由其承审员万某代理之时，一般县绅，以为千载一时之良机难得，乃纠集党羽，大开会议于某祠内，蜂拥县府，抢夺县印”。\footnote{《广昌社会调查》，《汗血月刊》第1卷第5期，1934年7月20日。}当时，地方官员如不能取得乡绅的背书，施政十分困难，甚至有遭到驱逐的可能。即使离南昌不远的东乡县也有乡绅因反对新任区长而将支持新区长的乡绅打伤，致该区长“难上任”。\footnote{《训令东乡县长据该县第三区长陈汝楫呈为抗令不交纠众胁迫恳请令县负责维持俾便到差视事》，《江西民政季刊》第1期，1930年1月。}为保持自身权位，大部分地方官员到任后都选择和乡绅妥协，或则各守分际，或则沆瀣一气，时人记载：“民国以来，任上犹县长者，虽不敢谓尽属坏人，但一经与土豪劣绅结合，即不坏者亦坏。”\footnote{《上犹人之赤祸谈》，《共匪祸赣实录》第2期，江西旅沪同乡会编印，1931，第57页。}虽然这些对绅权力量的描述不一定全面反映绅权的真实运行状况，但其负面效应的存在应属无疑，南昌行营曾判断：“一二匪党，不逞之徒，因缘时会，煽其共产邪说，诱惑裹胁，成今日燎原之势，推穷其故，岂民众生性好乱，毋以土豪劣绅之厉耳。”\footnote{《国民政府军事委员会委员长南昌行营政字第九四六号训令》，《湖北地方政务研究周刊》第1卷第12期。}


南京国民政府成立后，地方治理秉承孙中山的地方自治原则，这与积极的行政控制在理论上有所抵牾。而且，无论是传统中国既有体制还是当时对中国有着深刻影响的西方社会政治理论，都不倾向过于严密的行政控制，这对南京政府不能不有所影响，所以，在南京政府内部，自治、均权和民主的力量始终不可小觑。同时，国民党作为新兴的以革命自命的政治力量，显示力量、舍我其谁的集权呼声高涨，苏俄体制乃至德意国家的集权发展模式，对蒋介石及南京政府也形成相当刺激。由于此，南京政府建立后的相当一段时间内，来回拉扯于两种趋势之间，政治定位并不十分清晰。一面是经历十数年分裂状态的国家，一面是缺乏有效凝聚力的社会，国民党很长一段时间内都未在地方自治和中央集权之间找到平衡点。从江西、福建两省看，共同的特点是行政控制软弱，被赋予地方自治功能的县级行政机构一般只设五个局、科，公务人员很少。1930年代的统计，江西统计的43县有1237名公务员，平均每县不足30人；福建统计的45县只有公务员1134人，平均每县只有25人。江西44县中，县长兼办司法者36县，不兼办者8县；福建45县中，县长兼办司法者33县，不兼办者12县。\footnote{《县政调查统计·江西省》，国民政府内政部：《内政调查统计表》第22期，1935年6月；《县政调查统计·福建省》，《内政调查统计表》第21期，1935年5月。}大部分县份司法和行政混淆，体制之有待健全毋须多言。


由于政治形势混乱及控制不足，当时江西、福建两省县长平均任期不到一年，江西1930～1934年间更动县长265人，平均任期328日，最短的只有9天。\footnote{《县政调查统计·江西省》，《内政调查统计表》第22期，1935年6月。}福建同期更动县长289人，平均任期284日，最短的只有4天。\footnote{《县政调查统计·福建省》，《内政调查统计表》第21期，1935年5月。}多数县长任期在半年至一年间。这样频繁地更动地方官，使其无法真正了解、管理地方事务，地方权力更多落入士绅和胥吏手中，这对于地方政治和社会进步造成十分不利的影响。所以，南京中央政府成立后，在江西、福建两省几乎看不到地方政治的真正变化。


中共在赣南、闽西建立根据地后，其强大的政治、宣传组织动员力量迅速显现，地方政治结构在军事力量的支持下以阶级原则进行重组，地方豪绅很快遭到毁灭性打击。中央苏区如毛泽东所言，是枪杆子底下出政权，大多数政权都经由红军占领而建立。通常是先建立革命委员会，待政权巩固以后，再成立苏维埃。如下的报告基本可以代表红军到达后政权推进的基本路数：

\begin{quote}

地方工作现建立了五个支部，开了一个群众大会，选举了区革命委员会，为罗坊区，共有六个乡政府，分配土地，没收豪绅地主财产，现进行调查与分配工作，游击队的组织现还未弄好，不过受着我们之宣传，今已开始自动性了。工作团留了几个同志在罗坊负责，近日有可能的组织好，才能够巩固罗坊一带政权，使我们向前发展，减少一切困难。\footnote{《南丰中心县委给总政治部信——整顿党团组织与地方群众工作问题（1932年9月15日）》，《闽赣苏区文件资料选编》，福建三明档案馆、福建建阳档案馆、江西抚州档案馆编印，1983，第1页。}
\end{quote}


因此，中央苏区一开始就具有强烈的军事推进特征，这种军事力量对政治和社会的主导作用，是观察苏维埃运动须臾不可忘却的基本背景。即使到1934年苏区已经历空前发展时期后，张闻天仍然承认：

\begin{quote}

我们已经不止一次听到这样的话，说我们应该反对依靠大红军的观点，似乎大红军一走就什么工作冒〔冇〕办法的观点。但是这种反对，只有当我们真正能组织群众的武装斗争时才有意义，没有大红军，又没有群众的武装组织的依靠，而想在刀团匪骚扰的区域，建立一个空洞的苏维埃机关，这种机关显然经不起刀团匪的一次截击就会完全坍台的。\footnote{《闽赣党目前的中心任务——张闻天同志在闽赣战地委员会扩大会上政治报告的最后一节（1934年7月26日）》，《斗争》第71期，1934年9月7日。}
\end{quote}


由于红军的绝对强势地位，初期苏维埃政权本身反而相对较为薄弱。来自中共内部的报告指出：“许多苏维埃政权，不但不能为群众谋利益还要压迫群众，群众不认识苏维埃是自己的政权，不敢批评政府监督政府，所以这个政权还是脱离群众。”\footnote{《中央苏维埃区域报告》，《中央革命根据地史料选编》（上），第377页。}“因为苏维埃负责人脱离群众命令群众，在敌进攻时只顾自己逃命或瓜分公家金钱，群众没有分，不管群众的生死，所以永阳乡苏要伙夫也要捉来，全区民众无论一个什么行动，十分之九要捉群众，会议也不到，对红白军的认识无甚分别。”\footnote{《张启龙关于赣西工作情形给中央局的报告（1931年7月12日）》，《江西革命历史文件汇集（1931年）》，第109页。}一些仓促成立的政权领导不力、干部不足、组织经验缺乏、政府机关人员缺乏为群众服务的意识，因此组织不正规、工作中的随意性等问题尚普遍存在。当时报章记载：

\begin{quote}

从康都到黎川，有43个挑夫，走到中途，有些哭起来了。后来调查结果，内中有一个年纪已68岁，有一个63岁，一个61岁、56岁，到60岁的有4个，14岁的小孩，也有一个。问他们大家回答说：“我们洛口清太乡苏维埃，总是压迫人，派公事不公平，壮年人偏不派，硬要派我们这些老人，乡苏罗主席、赖秘书，总是打人。他们威风很大，拿到政府内就打屁股！还有一个姓丘的现当赤卫队长，他过去当过兵，非常恶，专门打人……罗主席穿得很好，专偷别人的老婆。”此外水南也发生过这样的事情，就是有一个挑夫因为不肯走，就要把他拿去枪毙。\footnote{《反革命分子，躲在苏维埃机关里捣乱》，《苏维埃》（中革军委机关报）第2期，1931年7月7日。}
\end{quote}


从当时材料反映的事实看，这样的事例并不罕见：“大多数的政府，群众并不认识为他们自己的政府，认为不过如象反动统治过去设的什么一样，有什么纠纷即到政府里去解决。有些群众实际叫政府为局，如我经过安福的桑田，我问农民政府设在什么地方，他不懂，结果他答复我是‘局’设在某处才带我去”。\footnote{《赣西南刘作抚同志（给中央的综合性）报告》，《中央革命根据地史料选编》（上），第247页。}


政府的软弱状况，还缘于当时军事紧张的实际形势。苏区建设初期，红军领导人毛泽东、朱德、彭德怀等主要精力理所当然放在军事上，而红军主力又不断处于运动开辟、保卫根据地的状态，难以在一个固定地区长期指导地方工作，这使主要依靠红军开辟、发展、巩固根据地的苏维埃政权基础先天不足。共产国际远东局曾引述江西苏维埃政府主席胡海1930年的说法：“该同志动身来上海（9月29日）之前，我们红军的指挥人员（毛泽东、彭德怀）与政府没有任何联系。政府是一回事，军队是另一回事。毛泽东对政府工作几乎不感兴趣。”由于此时毛泽东领导的总前委和赣西南地方政权在许多问题上存在分歧，不排除这一报告有刻意夸大的成分，但其反映的军队和地方沟通不畅的问题当非捏造。因此，共产国际远东局明确建议中共中央：“应该使毛泽东不仅对军队的状况和行动负有责任，而且还要让他参加政府并对政府的工作负有部分责任。应该任命他为政府委员（革命军事委员会主席）。”\footnote{《共产国际执行委员会远东局给中共中央政治局的信》，《共产国际、联共（布）与中国革命档案资料丛书》第9册，第451～452页。}


随着中共在赣南、闽西控制的逐渐稳固，苏维埃政权建设开始趋于健全，但政权在整个控制体系内的地位，仍然和传统政治体制有别。贯穿中央苏区之始终，苏维埃政权和军队、党的关系大体处于这样的框架中：军队是基础、政党是灵魂、政权是手足。


作为中共的第一次建政尝试，苏维埃政权十分强调人民性和民主监督。在肯定阶级分化的基础上，人民民主建设一直是中共孜孜以求的目标。武装革命早期，中共的这种民主理念主要体现在部队建设上。中共建军初期实行的士兵委员会制度，虽然后来由于与党领导军的原则有所冲突未能继续下去，但在当时中共组织力量有限的状况下，这一高度的军内民主制度对有效地凝聚部队力量，改善官兵关系，提高全军的责任心和战斗力发挥了不可低估的作用。


苏维埃政权成立后，中共着力建设一个人民性政权。苏维埃政权沿袭苏俄的体制和经验，实行代表会议制度，中华苏维埃工农兵代表大会及其中央执行委员会是中华苏维埃共和国的最高政权机关。中央人民委员会为国家行政机关，负责处理国家的日常政务。苏区居民凡年满18岁者，都有权“直接派代表参加各级工农兵苏维埃的大会，讨论和决定一切国家的地方的政治任务”，对苏维埃政权机关工作人员享有选举、监督、罢免和撤换权。


中华苏维埃共和国的地方政权，实行“议行合一”制。苏维埃政府基层政权分县、区、乡三级，通过选举产生。省、县、区苏维埃政府，均设立与中央政权相应的人员和办事机构。乡苏维埃政府机构较简单，但一般都设有各种经常和临时的委员会，如扩红、土地委员会等，在一些工作先进的乡，如长冈乡的常设委员会多达12个，包括扩大红军、土地、山林、建设、水利、桥梁、国有财产、仓库保管、没收、教育、卫生、防空防毒委员会等，这些委员会有许多一直设立到村，乡委员会的委员多由各村相关委员会主任代表组成，各村的工作可通过各委员会得到有效沟通。对此，毛泽东十分赞赏，认为：

\begin{quote}

乡的中心在村，故村的组织与领导成为极应注意的问题。将乡的全境划分为若干村，依靠于民众自己的乡苏代表及村的委员会与民众团体在村的坚强的领导，使全村民众像网一样组织于苏维埃之下，去执行苏维埃的一切工作任务，这是苏维埃制度优胜于历史上一切政治制度的最明显的一个地方。\footnote{毛泽东：《才溪乡调查》，《毛泽东农村调查文集》，第336页。}
\end{quote}


村代表制度使苏维埃政权的触角有效深入到村一级，虽然传统的保甲制度也具有这一功能，但村代表制度网罗的人员、领域、层级都远较保甲宽广，加上中共基层党组织的领导，使政权效能发挥相当充分：“他的命令，是如水之流下，由区乡政府，而党员，而直达群众。党员向群众传达命令的方式，对农则在田野，对工则在工厂，对商则在市廛。”\footnote{孔荷宠：《消灭土匪的意见》，南昌《市政半月刊》第1卷第5、6期合刊，1934年10月1日。}毛泽东所谈到的这样一种状况，在瑞金武阳区石水乡春耕动员中有十分具体的体现：

\begin{quote}

（1）动员的开始：乡苏主席参加区苏讨论春耕的会议回去以后，连同区苏、区委派去的指导员，首先开了党团员大会，随即开了乡苏代表会，接着开了贫农团大会，最后开了妇女代表会，赤卫军、少先队、儿童团下操时，都着重讲说了春耕。乡苏和妇女代表会都组织了宣传队，男女两宣传队均有宣传员三人（三村每村一人），队长一人，都作了春耕的宣传，贴了标语。这样就在全乡开始造成了春耕的热烈空气，大家晓得春耕的意义、目的、计划与进行的方法了。

（2）乡代表会的领导：乡代表会起了极大的领导作用。71个代表（内有妇女13人）分在三个村中，每人对其选民60人（工人、雇农的少一些）负起春耕领导之责，每十天乡代表会开会，报告各村状况，讨论新的办法。

（3）贫农团的领导：这里的贫农团有500多人，共有52个小组，每组8人至12人不等，分在三村，每村一个干事会。全乡一个总的干事会。领导春耕的方式：每五天各小组开会一次，每十天三村贫农团各自开会一次，每月开全乡贫农团大会一次。

（4）妇女代表会的领导：全乡96个妇女代表，选举主席团7人。每5天由各个代表召集所属妇女开会一次，每10天全乡妇女代表会一次。本乡妇女（全区亦然）在生产中占了极重要的地位，她们除了犁田、耙田、莳田之外，什么事都会做，铲革皮，割卤萁，开塘泥，开粪下田，作田塍、田坎，修陂圳，犁田时散粪砍菜子，耘田、巡水，种杂粮、蔬菜，样样都会。她们的劳动占了全乡生产劳动的50%以上，在武阳区犁田、耙田、莳田也有少数妇女会做，她们正在学习，使多数人都会做。

（5）儿童在春耕中：由于儿童团的动员，他们在铲草皮、看牛子、捡狗粪。他们规定三个月内每人捡100斤狗粪给红军公田，100斤给红军家属，100斤给自己，捡来了过秤，载在他们的比赛条约。

（6）党对春耕：党在春耕中起了极强的总的领导作用。全乡有党员150多，团员70多。1个总支（干事会11人），三村每村1个分支（干事会3人至5人）。为了春耕，5天开一次小组会，半个月一次分支会，一月开一次总支大会。

（7）耕田队：男子耕田队全乡1个大队，有大队长，队员520多人。分为4个分队，第一村人口800多，1个耕田分队，第二村人口1200多，2个耕田分队，第三村人口700多，1个耕田分队，每分队队员100多人，有分队长，下分小组，每组5人至10多人不等。女子耕田队400多人，组织与男子队同。男女耕田队共约千人，是专为耕种红军公田与帮助红军家属、苏维埃工作人员耕田而组织的，他们吃了自己的饭去耕，耕的特别好。

（8）突击队：党与政府各1个3个人的突击队，3个人分在三村，公开的号召，秘密的考查，各种春耕的会议他们都去参加，可是群众不晓得他们是突击队。

（9）比赛：全乡都发动了生产比赛，激发了群众极大的革命热忱。本乡的办法：家与家比，屋（小村）与屋比，村与村比，乡与乡比，比什么？比耕耘、比肥料、比杂粮蔬菜种得好，党、团、政府与群众团体会议的紧张，工作的激进，比赛起了大的推动作用。\footnote{《中央土地人民委员部关于夏耕运动大纲（1933年4月22日）》，《革命根据地经济史料选编》（上），江西人民出版社，1986，第256～257页。}
\end{quote}


从上面的材料可以看到苏区组织的严密体系：党、政府、群众组织，环环相扣，层层相连，整个动员过程有宣传、有组织、有监督、有领导，堪称天衣无缝。不过，多少让人觉得有些突兀的是，春耕本为农民的天然义务，中共为何要在这方面付出如此多的心力，这可能和一段时间内中共土地政策不稳定，农民生产积极性受到挫伤有关。从长远看，解决这些问题，更需要的应该还是政策调整，以利调动农民的积极性，而不仅仅是大规模的组织动员，因为最有效的动员，也有日久生疲的可能。当年的材料显示，苏区动员很少会谈及或考虑个人需要，这有其特定的理论背景，但无论当时还是日后，利益调整都是必须面对的现实，忽略了这一点，许多的政策制定可能就会有空中楼阁的风险，毕竟，把理想等同于现实未必就是实现理想的捷径。


尽管不无可讨论之处，苏维埃政权的群众动员总体上看还是反映着依靠群众、发动群众的人民性立场。苏维埃制度理论上是建构在民主运作的框架之中。1933年底召开的江西省第二次工农兵代表大会明确要求：“对城乡苏维埃的工作在广泛运用民主主义的基础上，要实行选民的召回权，建立村代表会议及村代表主任制，在乡和村中组织各种委员会，吸引广大群众参加苏维埃工作。”\footnote{《江西省第二次工农兵代表大会决议案（1933年12月28日）》，《江西革命历史文件汇集（1933～1934年）》，第333页。}选民召回权的强调，使代表的权力和责任与选民发生直接的联系；而委员会制度既保证了上下通联，又体现了群众意志。1933年5月，安远县苏维埃政府要求各地设立茶亭，方便来往行人，对于设立茶亭的经费，县府要求“当地政府负责召集当地乡苏代表会议和各群众团体会议，讨论决定其费的开支，在（再）由当地群众募捐供给之，切不可借口滥捐”。\footnote{《安远县苏维埃政府内务部训令》，《江西革命历史文件汇集（1933～1934年）》，第128页。}可见在当时的政权运作体系中，如果充分发挥制度设计的效能，代表会议可以成为反映民意和决定政策的重要一环。中共在实践中摸索出的这一基层政治运作体制，虽然由于苏区的现实环境、党组织与乡村政权的磨合需要及时局变化等因素未及真正展现和推广，但体现了中国传统乡村政治面貌的充分潜力。


为加强各级苏维埃政权的建设，临时中央政府先后制定与实施了苏维埃选举细则，对苏维埃的选举原则及方法作出规定，并领导各级苏维埃展开由下而上的选举运动。中共有关决议明确指出：“苏维埃的民主，首先表现于选举运动。”\footnote{《江西省第二次工农兵代表大会决议案（1933年12月28日）》，《江西革命历史文件汇集（1933～1934年）》，第355页。}苏维埃代表选举在当时受到相当的重视。从选举过程看，虽然中共党组织要求不要直接干预选举结果，强调“党对于苏维埃委员选举可以拟定名单，但不能直接向群众提出，只有在群众中宣传，使群众接受这种意见，选举出来”，\footnote{《目前政治形势与党的组织任务（1930年7月22日）》，《中共中央文件选集》第6册，第216页。}但包办的现象特别在初期事实上难以避免。随着中共苏区政权的巩固，选举也逐渐制度化，根据中共中央的要求，1933年进行的苏维埃代表选举中，各地注意“在选民大会中必须发动选民斗争，从斗争中来选举工作积极、斗争坚决的真正劳苦工农青年及劳动妇女充当苏维埃代表，反对过去指派方式”。\footnote{《安远县苏维埃政府关于区乡苏维埃改选运动的训令（1933年4月9日）》，《江西革命历史文件汇集（1933～1934年）》，第94页。}赣县的选举“相当的发动了群众的斗争，如在选举时选了不好的分子当代表，群众立即反对（清溪古村区）”。\footnote{《江西赣县两个月（10.20～12.20）冲锋工作报告（1933年1月10日）》，《江西革命历史文件汇集（1933～1934年）》，第6页。}兴国、上杭、瑞金等地“普通的以乡为单位，组织了三人至七人的宣传队，比较先进的地方组织了化装讲演，演新戏，俱乐部开晚会，各学校上选举课等，造成选举的热潮”。\footnote{梁柏台：《今年选举的初步总结》，《红色中华》第139期，1934年1月1日。}根据毛泽东1933年的调查，这年年底进行的苏维埃代表选举虽然没有实行直选，而是层级代表选举，但有些地区选举还是开展得有声有色。福建上杭才溪乡公布的候选名单得到群众积极回应：

\begin{quote}

下才溪一百六十多人（内应选九十一人），一村贴一张，每张均写一百六十多个名字。群众在各人名下注意见的很多，注两个字的，五六个字的，十多个字的，儿童们也在注。注“好”、“不好”等字的多，注“同意”或“消极”的也有。有一人名下注着“官僚”二字。受墙报批评的有二十多人，被批评的都是只知找自己生活、不顾群众利益、工作表现消极的。有诗歌。内有三张批评乡苏对纸业问题解决得不好。\footnote{毛泽东：《才溪乡调查》，《毛泽东农村调查文集》，第337页。}
\end{quote}


同样的状况在各地都有显现，兆征县大埔区教育部长“在全区工农兵代表大会上，要他来把全区文化教育工作，向代表作报告，结果他讲了七句话。因之，全区的代表同志，很热烈的站在无产阶级的立场向他作了无情的斗争”。\footnote{《打击文化战线上的官僚主义者》，《红色中华》第142期，1934年1月10日。}胜利县古龙岗中邦乡“群众对候选名单上的贪污腐化分子、开小差分子只在夜晚在公布的候选名单上的姓名底下写字”。\footnote{《中共江西省委关于全省选举运动的检阅的通讯（1933年10月）》，《江西革命历史文件汇集（1933～1934年）》，第288页。}选民的批评热情体现了其对选举的主动参与和重视。


民众在选举运动中用选票真正体现自己的意志，可以鼓励民众对政治的参与意识，加深对政权的信任度与休戚感，并反映出民意的趋向。1933年底的苏维埃代表选举，选民的参与程度相当高，“先进的如：兴国全县，上杭才溪区、瑞京武阳区平均到会的选民都占百分之九十以上。比较落后的地方如西江县洛口县，到会的选民平均在百分之六十二以上。中等区瑞京的下肖区到会选民平均在百分之七十一以上”。选民在大会上提出的提案包括诸如“扩大红军，优待红军家属，解决几年来无音讯的红军老婆的问题，消除市面现洋与纸币的差异现象，进行节省运动，推销公债，扩大合作社组织，准备春荒，女子由甲地嫁到乙地土地问题解决，修理道路桥梁，设俱乐部列宁小学等”。这些提案集中在与民生密切相关的具体问题上，相当程度上“反映出群众的要求”。\footnote{梁柏台：《今年选举的初步总结》，《红色中华》第139期，1934年1月1日。}同时，群众还可通过选举批评监督和罢免干部，对不称职的干部，在选票名字下面注“消极”、“不好”、“官僚”等字样。对犯错误的干部，只要选民十人以上提议、半数以上选民同意，就可予以罢免。


对于一场中国历史上崭新的基层民主实验，考虑到当时历史条件、环境和当时中共建政思想的限制，各方面的民主理念和民主实践都严重不足，过高估计其成果也是不切实际的。各地比较普遍的反映是选举中缺乏对工作的充分批评，候选人与选民之间的交流也远远不够：“在各村选民大会上，乡苏主席都向选民作了工作报告，但是不很充分，所以不能发动群众起来热烈批评政府的工作，这是一个大的缺点。”\footnote{《博生县梅江区的选举运动》，《红色中华》第122期，1933年10月27日。}即使是毛泽东表扬的模范乡的长冈乡也存在下列缺点：

\begin{quote}

（1）宣传没有指出，苏维埃是群众自己管理自己生活的政权，选举苏维埃代表是群众最重要的权利。（2）候选名单人数恰如应选人数，没有比应选人数增加一倍，因此群众对于候选名单没有批评。选举委员会在组织候选名单问题上没有起什么作用，只有党的活动。（3）工作报告会议上没有尽力发动群众对乡苏工作的批评。\footnote{毛泽东：《长冈乡调查》，《毛泽东农村调查文集》，第297页。}
\end{quote}


在战争环境下，有的地方问题可能还更加严重一些。边区县区常反映：“候选名单没有充分准备好及充分宣布充实区代表权利，与被选为代表的任务和光荣，以致有的竟临时来拉，所以部分当选人，不愿意当代表。”\footnote{《万崇区四乡选举运动》，《红色中华》第122期，1933年10月27日。}赣县“田村乡的选举不严密以致流氓混入了来开选举会，边区的有些乡在选举时群众怕来开选举会，怕会当代表”。\footnote{《江西赣县两个月（10.20～12.20）冲锋工作报告（1933年1月10日）》，《江西革命历史文件汇集（1933～1934年）》，第6页。}对于缺乏民主政治经验的人们来说，要准确理解选举和代表大会制度就有困难：

\begin{quote}

群众对于苏维埃的认识很微弱，只知道埃政府是他们的政府而不知道“埃政权”的内容，一班农民都应有实际参加政权的权利，就是常常要召开群众大会代表大会讨论他们自己的事情。他们并不知道政权有二，一为代表会议，一为执委会，不知道代表会议才是合法的政权机关，而执委会不过是代表会闭会后，受代表会委托的执行代表会议的处理日常事务的机关，更不知在一哄而聚的群众大会之中能讨论问题，必须假手于代表会议。他们以为选出几个人坐在机关里，就叫做苏维埃，所以各级机关时常只有执委会议（甚至执委会还少开，只有主席、财务、赤卫或秘书只几个人处理一切）而没有代表会议。这是一般群众政权的意识大薄弱，而一般同志对于政权的认识也是莫名其妙的缘故。\footnote{《杨克敏关于湘赣边苏区情况的综合报告（1929年2月25日）》，《中央革命根据地史料选编》（上），第49页。}
\end{quote}


在此基础上，选举委员会出现组织不合法的状况也就不难想象：广昌长桥区被监督者成了选举的监督者，“选委会只由区苏主席、各部长组织，以主席为主任”。\footnote{《中共江西省委关于全省选举运动的检阅的通讯（1933年10月）》，《江西革命历史文件汇集（1933～1934年）》，第288页。}熟悉代表会议制度的人们对这种角色错位可能会啼笑皆非，但当年的操作者却认为理所当然。苏维埃代表选举在开了赣南、闽西地区群众意志表达方式之先河的同时，其出现的种种问题，却也不可避免影响到选举的效果，使选举的意义有所遮蔽。

\subsection{苏维埃政权的监督体系}


由于主、客观条件的制约，选举尚难成为苏维埃政权自我监督、自我更新的足够有效方式，代表会议制度也往往流于形式：

\begin{quote}

乡、市苏的工作，大多数都是主席独裁与包办。乡、市苏的代表，除了照例参加会议外，很少切实过问乡、市苏的工作；乡、市苏代表，未能将自己所代表群众的意见和要求，反映到乡、市苏来；乡、市苏的决议，亦很少带回去向选举自己的选民作报告，动员群众来执行。一般乡、市苏代表，尚未能成为群众斗争与生活改善强有力的组织者与领导者。\footnote{《省苏执委会工作报告的决议案（1933年3月23日）》，《闽浙赣革命根据地财政经济史料选编》，厦门大学出版社，1988，第156页。}
\end{quote}


在此状况下，中共为维持政权的有效运作，更多的还是依赖严格的纪律和组织、监督体系。


1932年后，中央苏区苏维埃政权建立了一套组织严密与面向群众相结合的监督机制。各级党组织设立监察委员会，各级政府设立工农检察部和工农控告局，负责对干部进行监督。1934年2月后，中央审计委员会对中央政府各直属机关、企事业单位和群众团体进行普遍稽查，对财政开支状况实施审核。


发挥群众的监督作用，既反映着苏维埃政权的人民性，也是廉洁政治的重要一环。1932年9月中央苏区颁布《工农检察部控告局组织纲要》，规定“在工农集中的地方”，可设立控告箱，群众随时可以控告、揭发苏维埃政府工作人员的“贪污浪费、官僚腐化或消极怠工现象”，还可“指定不脱离生产的可靠工农分子，代替控告局接收各种控告”。\footnote{《工农检察部控告局组织纲要》，《中华苏维埃共和国法律文件选编》，江西人民出版社，1984，第417页。}当时规定必须是实名控告。不过，实际执行中并未画地为牢，董必武任苏维埃最高法院院长期间，曾处理控告中央总务厅采买员的匿名信，发现不尽属实后，查找到写匿名信者，未按诬告处理，而是加以劝导、说服。中共党人严格的监督规定和细致耐心的工作，保护和激发了群众控告的积极性，当时查处的案件不少来自民众的直接控告。


为保证群众有效参与对干部的监督，各级工农检察部经常组织“轻骑队”、“突击队”等，对各级部门进行检查、监督。轻骑队是青年团领导下的青年群众组织，最早成立于1932年7月，主要任务是：“检查苏维埃机关内、企业内、经济的和合作社的组织内的官僚主义、贪污、浪费、腐化、消极怠工等现象，举发对于党和政府的正确政策执行的阻碍与曲解。”\footnote{《轻骑队的组织与工作大纲》，《斗争》第41期，1934年1月5日。}突击队苏区内凡有选举权者均可加入，归当地工农检察部领导，可对苏维埃机关和国家企业进行突击检查。轻骑队和突击队通过微服私访，明察暗查，发现问题及时报告有关部门查处。


苏区对干部的贪污、浪费行为惩处极严。1930年3月，闽西第一次工农兵代表大会通过的《政府工作人员惩办条例》，规定下列行为即可予以枪决：“侵吞公款至三百元以上者，受贿至五十元以上者。”\footnote{《政府工作人员惩办条例》，《中央革命根据地史料选编》（下），第87页。}1933年12月15日，苏维埃中央政府发出训令，规定：凡“贪污公款在五百元以上者，处以死刑”；“贪污公款在三百元以上五百元以下者，处以二年以上五年以下的监禁”；“贪污公款在一百元以上三百元以下者，处以半年以上二年以下的监禁”；“贪污公款在一百元以下者，处以半年以下的强迫劳动”；以上罪犯同时“没收其本人家产之全部或一部，并追回其贪污之公款”。训令同时宣布：“凡挪用公款为私人营利者以贪污论罪”；“因玩忽职务而浪费公款，致使国家受到损失者，依其浪费程度处以警告、撤销职务以至一个月以上三年以下的监禁”。\footnote{《中央执行委员会廿六号训令》，《红色中华》第140期，1934年1月4日。}据此，中央苏区不少干部遭到严肃处理。1933年4月，江西省苏维埃政府以伙同他人盗卖鸦片罪判处原胜利县苏维埃政府主席钟铁青死刑，原胜利县临时县委书记钟圣谅监禁两年。1933年5～8月，广昌因贪污问题被清除出党的党员干部就有6人，包括区委书记、区苏主席、区苏军事科长等。\footnote{《党的组织状况》，《中央革命根据地史料选编》（上），第695页。}安远龙布区苏维埃主席团会议记录显示，该区1933年3月工农检查部发现主要问题有：

\begin{quote}

（一）长河乡财政贪污浪费，文书腐化；（二）上林乡主席朱文求纠结红军家属老婆；（三）增坪乡古山消极怠工不能打理工作……（五）各乡通知文件都忽视了；（六）正副主席正主席余恭梅包庇富农分子，副主席李荣德没有一点打理工作，几件事军阀手段苏维埃压迫来负责。对于成立合作，自己合股又分三家布X；（七）特务排长叶桂芳在前方开小差回家，班长张天祥开小差回家，三班长谢新荣独立团开小差。\footnote{《安远县龙布区苏主席团会议记录（1933.3.10～6.16）》，《江西革命历史文件汇集（1933～1934年）》，第48页。}
\end{quote}


可以说，苏维埃工作人员事无巨细，都处在政权机关的严密监督之下。


中共以严刑峻法和严密监督努力抑制贪污腐败，同时通过对社会的有效控制扫除社会长期累积的腐败因子，确实取得了很好的效果，苏区吏治的廉洁一直是人们津津乐道的话题，被认为是“空前的真正的廉洁政府”。\footnote{《关于四个月节省运动的总结》，《红色中华》第232期，1934年9月11日。}当时的苏区干部可以从政府方面获得的利益主要就是苏维埃政权规定的：“政府工作人员，在服务期间，本人及其家属（即父母妻子）可按照应纳之国税额减收半数，用以减轻政府工作人员的家庭经济的负担。”\footnote{《福建省苏维埃政府通令第二十五号》，《福建革命历史文件汇集》甲15册，第304页。稍后又指出：“必须是政权机关的工作人员，才可以减税，其他的机关团体，如共产党、青年团、工会、互济会、少先队、反帝大同盟……这都不是政权机关，而是革命群众团体的组织，不能与政权机关相提并论，其工作人员不能减税。”（《福建省苏维埃政府通知第二十号》，《福建革命历史文件汇集》甲15册，第321页）1933年颁布的《农业税暂行税则》也规定，政权机关工作人员及其父母、妻子可减征50\%的农业税。}但是，他们在工作上投入的精力显然要多于所得，当时干部回忆：

\begin{quote}

我们去县里开会时，每人自备一个席草饭包，吃多吃少由自己定。在用膳前，用自己的饭包装自己带去的米或薯米，然后把饭包口扎紧，做上记号，放在一口锅里统一煮……如果县、区、乡召开只有一天的会议，来开会的多数人，一不带饭，二不带干粮，勒紧裤带，饿着肚子，一直坚持开到底。

杨尚奎同志任村文书时，经常早饭后便去村苏维埃政府工作，一直到天黑才回家。中午这餐不是吃干粮，就是勒紧裤带饿肚子，晚上，有时回得太晚，家里留的稀饭冷了，就吃冷稀饭，如果没有留什么吃的，有红薯，就吃几只生红薯充数，没有红薯，只好喝碗冷茶或冷水，填一填肚子。\footnote{万香：《回忆苏区干部作风片段》，《兴国文史资料》第1辑，政协兴国文史资料工作组编印，1982，第79～81页。}
\end{quote}


当然，腐败问题的解决不可能一劳永逸，苏维埃政权也并不是白璧无瑕。在当时战争环境下，权力的运用经常要超出体制化的监督，在基层尤其如此。民众对干部自觉的监督意识不足，选举制度在当时情况下更多的还是走过场，而且选举仅限于苏维埃代表，政府人员不必直接面对民众的选择。监督渠道从制度上虽多有开辟，但许多地区也不尽畅通，如有些政府设立控告箱时走过场，安远县龙布区“全区成立控告箱四只”，\footnote{《安远县龙布区苏主席团会议记录（1933.3.10～6.16）》，《江西革命历史文件汇集（1933～1934年）》，第48页。}平均一乡还不到一个，其所发挥的作用自然大受限制。因此，和对传统社会的腐败现象扫除取得的巨大成就相比，政治权力控制相对要弱一些，贪污违纪现象在苏维埃区域内仍然存在，某些地区、部门甚至相当严重。湘鄂赣的酃县“三区乡负责者，每人分得好田十一担谷子，群众有四个一家只分得十一担谷子的差田……群众要责问，马上‘反抗政府’‘反动派’的帽子即戴上头来了，有许多群众反对政府负责人的官僚腐化分子，但又不敢公开说”。\footnote{《中央局滕代远巡视湘鄂赣苏区的报告（1931年7月12日）》，《湘鄂赣革命根据地文献资料》第1辑，第524页。}


当时，腐败最严重、最具有警示意义的还是“于都事件”。1934年初，中共中央在检查工作时，发现于都各级干部贪污严重。随后，中共中央连续组成调查组到于都，项英也亲往于都督促调查，在此发现“四十件各种各色的贪污案件”，涉及从县到乡的各级机关和大批干部：“城市区苏三个主席九个部长，就有三个主席（区正副主席、工农检察委员会主席）六个部长（土地、劳动、内务、国民经济、财政、裁判）都是做投机生意的”；“贪污的种类也分几种。贪污分子由县主席部长以至乡代表。贪污成为风气，大家反不以为异，而且互相包庇、互相隐瞒”。\footnote{项英：《于都检举的情形与经过》，《红色中华》第168期，1934年3月29日。}对此，苏维埃中央政府作出严肃处理。中央政府执行委员、于都县苏主席熊仙璧，以权谋私，挪用公款做私人生意，于1934年3月被逮捕，撤销一切职务，判处监禁一年。于都县委书记刘洪清包庇熊仙璧的错误，带头利用职权，拉股经商，谋取私利，被撤销职务。县苏军事部长刘仕祥等冒领、贪污公款410元，并贪污打土豪缴获款项做私人生意，被判处死刑。同时以贪污罪被判死刑的还有少共县委书记滕琼、潭头区苏财政部长等5人。


于都县干部整体腐败案暴露出苏维埃政权的干部机制仍然存在的一些问题，如张闻天所写到的：

\begin{quote}

县委书记刘洪清当粤赣省委决定撤消他的工作时，他在会议上照着党委的批评依样画葫芦的承认了一些错误，其它的许多严重问题他是一句也不说的。他对于其它许多同他一样犯错误的人，也是一句话不说的。其它的县委负责人，也同样的仅仅批评了刘洪清已被揭破了的错误了事，多一句也是不肯说的。在撤消县苏主席熊仙璧的主席团会议上，也同样的重复了这一幕滑稽剧。甚至粮食部副部长会把与这一会议完全不相干的，而且同他自己工作也不相干的统计，如红军家属多少，男女多少等等来念一顿，但对于当前的问题，是愿意表示完全的沉默。\footnote{张闻天：《于都事件的教训》，《斗争》第53期，1934年3月31日。}
\end{quote}


干部内部相互包庇，串通一气，缺乏原则立场，同时群众对干部还没有形成真正的有效监督，这就给干部腐败提供了温床。在于都问题暴露后不久，赣县又发现据称不比于都问题小的贪污腐败案，“赣县县委的问题与于都县委无二样，而且更加严重”。赣县“区一级机关中投机商人更可以顺利的提到领导机关来，利用国家机关盗窃苏维埃的财产，大做其投机生意，江口区委区苏及江口商店就是样本的例子”。该县财政部门负责人刘绍稷“占据苏维埃财政机关公开的贪污，把苏维埃所没收的鸦片烟蒸过后勾结富农去贩卖，可以拿成千的银钱。用地主豪绅房屋的材料，不分给群众而为自己造大栋房子，每日酒肉，过其腐化生活”。\footnote{《把反雷开平两面派机会主义的斗争开展到全省去——省委对赣县工作的决定》，中共江西省委《省委通讯》特刊，1934年5月10日。}可见于都的问题并非特例。


中央苏区的苏维埃政权创造了中共历史上值得书写的民主范本，包括党内的批评可以相当充分地开展，报纸和舆论的监督比较严厉，乡村的民主选举有一定操作意义，轻骑队等具有官民合作背景的组织活动积极等，这些都使中共政权呈现出了自古未有的活力，群众对政治的参与空前广泛。苏维埃政权通过政权内部的监督机制和民众、舆论的外在监督予权力腐败以严厉控制，为中共廉洁政治、革新社会提供了成功经验。虽然由于权力的腐蚀，腐败现象仍然难以根除，产生腐败的温床在新的社会结构下还会以另一种方式出现，但中共当年在有限时间和条件下作出的有益尝试，仍然值得后人认真总结、谦虚面对。

\section{宣传功能下的教育和文化}


中央苏区所在的江西、福建在全国都不属于教育发达地区，以江西为例，该省1930年代全省县教育经费尚不及江苏两个大县教育经费之和。1929年的数据显示，江西儿童入学率在全国居第26位，初等教育经费居第18位，每人负担初等教育经费比居第25位，均在全国平均线以下。\footnote{程柏庐：《从二十年度全省教育统计上所得之感想》，《江西教育旬刊》第7卷第2期，1933年9月11日。}赣南、闽西在两省中又都属于偏僻、贫穷地区，教育基础更差。


在中共的革命话语中，苏维埃革命不仅是一场武装革命、政治革命，同时还是一场思想革命、社会革命，因此，对群众宣传革命理念，改变苏区民众教育文化落后状况，是革命的题中应有之义。中共对宣传鼓动向来十分重视，曾任粤军香翰屏部参谋的李一之深入观察中共宣传活动后发现：“全部红军有整个的宣传计划，各部队红军皆能取一致之步调，同一之意思，划一之口号标语以为宣传。”“共匪所至，字迹不拘大小优劣，必在墙壁遍涂标语，或标贴文字宣传，五花八门，不一而足，至其对俘虏对民众等亦能尽量宣传”。\footnote{李一之：《剿共随军日记》，第二军政治训练处1932年印行，第102页。李一之一生经历饶有意味：1926年入黄埔军校第三期步科，曾参加第二次东征和北伐战争。后任香翰屏部参谋，香被贬黜后，返乡从事教育。1937年任广州防空司令部少校参谋。抗日战争爆发后，任第十二集团军防空指挥部参谋，1938年5月秘密加入中国共产党。1940年任第四战区东江游击指挥所参谋、作战科中校科长，东江游击挺进纵队基干大队大队长，第四战区战地党政委员会政工大队大队长。同年9月因暴露党员身份，入广东人民抗日游击队第五大队，任军事教官、指导员。同年12月在东莞横涌被汉奸杀害。}中共的宣传有的放矢，用心缜密，陈毅曾举例加以说明：“比如红军标语打倒土豪劣绅这样写的时候很少，因为太空洞而不具体，我们必需先调查当地某几个人是群众最恨的，调查以后则写标语时就要成为打倒土豪劣绅某某等，这个口号无论如何不浮泛，引起群众深的认识。”\footnote{《陈毅关于朱毛军的历史及其状况的报告（一）（1929年9月1日）》，《中共中央文件选集》第5册，第762页。}


标语口号简单直接，明白易懂，视听冲击强烈，对文化水平不高人群宣传效果好。曾被红军监禁的外国传教士回忆：“他们所到之处，大写标语，红的、白的、蓝的，一个个方块字格外醒目……这些标语都是由宣传班写成的，这些人走到哪儿总是带着一桶漆，凡是能写字的地方，显眼的地方，他们都写大标语。”\footnote{《外国人笔下的中国红军》，陕西人民出版社，1996，第238页。}国民党军占领后进入苏区的调查人员也发现，红军标语“满坑满谷，随处可见。而国军到时，则每于标语相当之处，涂改数字，以为国军之宣传焉”。\footnote{贺明缨：《江西省田赋清查处实习报告书》，萧铮主编《民国二十年代中国大陆土地问题资料》第170种，第85089页。}一个是书写不厌其烦，一个是涂抹都敷衍了事，国共两党宣传、办事用心程度的差异，于此真是昭然若揭。


在苏区内部，中共更将宣传和教育相结合，《中华苏维埃共和国宪法大纲》规定：中华苏维埃政权以保证工农劳苦民众有受教育的权利为目的，在所能做到的范围内，开始实施完全免费的普及教育；积极引导青年劳动群众参加政治和文化的革命生活，以发展新的社会力量。《中华苏维埃第一次全国工农兵代表大会宣言》宣示：一切工农劳苦群众及其子弟，有享受国家免费教育之权。教育之权归苏维埃掌管，“政权组织教育机关与宗教事业绝对分开”。


第五次反“围剿”开始前后，中共对文化教育的重视不减。1933年4月，中央教育部发布第一号训令《目前的教育任务》，指出苏区文化教育的任务，是用教育与学习的方法，启发群众的阶级觉悟，提高群众的文化水平与政治水平，打破旧思想旧习惯的传统，以深入思想斗争和阶级斗争，参加苏维埃各方面的建设。同年7月，中央教育部发布训令，强调：“在国内战争环境中，苏区文化教育不应是和平的建设事业，恰恰相反，文化教育应成为战争动员中一个不可少的力量，提高广大群众的政治文化水平，吸引广大群众积极参加一切战争动员工作，这是目前文化教育建设的战斗任务。”\footnote{《中华苏维埃共和国中央教育人民委员部训令第四号——文化教育工作在查田运动中的任务》，《中央苏区革命文化史料汇编》，江西人民出版社，1994，第53～54页。}1934年1月，毛泽东在第二次全国苏维埃代表大会上提出苏维埃文化教育建设的中心任务：厉行全部的义务教育，发展广泛的社会教育，努力扫除文盲，创造大批领导斗争的高级干部。


中共在教育方面的努力，颇令蒋介石折服，他不无心仪地谈道：

\begin{quote}

匪区里面最紧张的，就是教育！最有纪律的，就是教育！最有精神的，也就是教育！而我们现在各地方的情形却不然，比方崇仁地方，所有的高小学校就完全停下来了，土匪他们什么经费可以少，教育经费一定要筹到，我们却反而要常常拿教育经费来做旁的用。\footnote{蒋介石：《以自强的精神剿必亡的赤匪》，《先总统蒋公思想言论总集》第11卷，第130页。}
\end{quote}


不过，如果深入观察中共在教育方面的作为，可以发现，蒋这段话其实并未真正发现问题的关键，苏区在教育上的投入经费并不像其想象的那样多。湘赣苏区1932年9月至1933年8月一年中，共支出252612元，其中用于教育的经费为275元，只占总支出的1‰强。\footnote{赵可师：《赣西收复区各县考察记》（四），《江西教育旬刊》第10卷第8期，1934年8月1日。}就这一点而言，远远无法和同时期国民政府3%左右的教育支出相比。中共开展教育的主要思路是利用民间力量、民间财力办教育，学校经费“原则上或由学生自纳，或由人民捐款”。同时大力提倡因地制宜、节省办学：“黑板常利用祠堂、庙宇之牌匾，加以刨平，涂以光油；或即就壁上刷黑一块。灯油粉笔，由学生自备。”\footnote{赵可师：《赣西收复区各县考察记》（四），《江西教育旬刊》第10卷第8期，1934年8月1日。}


因为教育的目标更多是立足于现实的思想宣传，中共对教育普及的重视远超过正规的渐进教育。在当时背景下，加强基础教育是迅速提高苏区人民文化水平相对现实可行的道路，因此，小学教育最受中共重视，苏区实施的五年义务教育以推广小学教育为目标。1933年10月，中央文化教育建设大会特别对教育部门在小学教育上的失误提出批评，认为其“对于组织小学教员的问题上，没有系统的工作”，\footnote{《中央文化教育建设大会决议案》，《中央苏区革命文化史料汇编》，第67页。}要求进一步加强对小学教育的管理、发展。为稳定、发展小学教育，1934年2月16日临时中央政府颁布《小学教员优待条例》，明文规定小学教员的生活费与当地政府工作人员同等，并同样享受群众帮助耕田和减纳土地税、免费治病等待遇；城市没有分田的小学教员，可给以其他物质上的帮助或资助其家属；小学教员每年给奖一次，连续取得第一等奖金两次者，按年增加原有奖金十分之二至十分之三。通过苏维埃政权的努力，中央苏区小学教育发展相当迅速。1933年底，江西、福建、粤赣3省的2931个乡中，设有列宁小学3052所，学生89710人，学龄儿童的多数进入小学读书，兴国儿童入学率达到60%。\footnote{《中华苏维埃共和国中央执行委员会与人民委员会对第二次全国苏维埃代表大会的报告》，《中央革命根据地史料选编》（下），第329页。}红军离开苏区后，1935年，由国民政府福建教育厅组织的视察注意到：“劳动小学，长汀一县，有三百余所。在小学课程，更注意课外活动，其目的都是注重关于守哨、侦探、放哨等工作。”\footnote{徐泰咸：《视察闽西的感想》，《福建特教通讯》第1卷第1期，1935年9月1日。}


和重视小学教育的出发点一样，为提高苏区民众文化水平，加强宣传效果，在文盲占人口绝大多数的苏区开展社会教育势在必行。1933年底，苏区共设有夜校6462所，受教成人达94517人，尤其是“妇女群众要求教育的热烈，实为从来所未见”。\footnote{《中华苏维埃共和国中央执行委员会与人民委员会对第二次全国苏维埃代表大会的报告》，《中央革命根据地史料选编》（下），第330页。}扫除文盲成为大规模的群众运动，各乡均要求组织识字团，由文化委员负责，“编十人一小组，设组长一人，三日召集组长前来教以三二个字，再由组长教各团员，并写字绘图在黑板或大纸上布置在通衢大道”。\footnote{《△△区苏维埃政府会议记录》，《江西革命历史文件汇集（补遗部分）》，第483页。文件未标区名，当是不明所致。从文件中以湖雷为中心展开工作看，这份文件属于福建永定县湖雷区，文中还提到湖雷附近的坎市、溪口更可以证明这一点。}据兴国高兴圩黄岭乡报告，该乡设有教育委员会，“教育委员会之下有识字班——全乡九十多班，四五人组一班，十六岁到四十五岁的男女都加入”，通过努力，“现在七岁以上的男女都识得字”。\footnote{《一个模范支部的工作报告》，《斗争》第5期，1933年3月15日。}同时，各地广泛设立俱乐部、书报社，“各乡至少要组织一个，乐器向民间搜集及政府购办，大村庄尽可能订新报，其油火由政府供给”。\footnote{《△△区苏维埃政府会议记录》，《江西革命历史文件汇集（补遗部分）》，第483页。}在苏区工作中走在前面的兴国，“全县列宁小学有四百八十余所，夜校一千余所。经过这些文化机关，连续不断的提高着群众的政治文化水平，教育着广泛的文盲的劳苦工农，所以兴国群众对于革命的认识比较的更深刻，对于斗争的参加也比较的更勇敢更坚决了……现在党团员占全人口百分之一十以上，并且建立了强有力的代表制”。\footnote{江西十县参观团：《参观兴国以后的感想》，《红的江西》第18期，1933年1月29日。}根据国民党方面收复苏区后的调查：

\begin{quote}

每一伪乡政府之文化部，必到处设有识字牌，每日书有宣传性或麻醉性之语句强迫附近民众，无论男女老幼务必前往认识并通晓其语意而后止。故当该县初收复时，试任叩一儿童以“阶级斗争”、“无产阶级”或“资本主义”等之意义，彼必能不假思索，对答如流，一如素有研究者然。土匪亦往往设立学校（如列宁小学等），但设备简陋，有黑板一方（且有时即利匾额改制）即可开设，桌椅粉笔等，均由学生自备，教师仅稍识之无者，即可充任。授课时间无一定，辄利用为匪军服役之余暇以为之（如放哨侦探等，土匪时利用儿童担任），有时晚间也须上课。\footnote{赵可师：《赣西收复区各县考察记》（三），《江西教育旬刊》第10卷第4、5期合刊，1934年7月11日。}
\end{quote}


苏区教育注重通俗易懂，如当时小学一年级的语文教材为《三字经》，第一页就是：“天地间，人最灵。创造者，工农兵。男和女，都是人。一不平，大家鸣。”\footnote{李柞煌：《江西省立列宁师范学校》，《回忆中央苏区》，江西人民出版社，1981，第380页。}其非常简明地传递出阶级观念、男女平权和反抗意识。在湘赣苏区，“无论男女老幼，都能明白国际歌，少先歌……及各种革命歌曲，尤其是阶级意识的强，无论三岁小孩，八十老人，都痛恨地主阶级，打倒帝国主义，拥护苏维埃及拥护共产党的主张，几乎成了每个群众的口头禅”。\footnote{《赣西南特委刘士奇给中央的综合报告》，《中国工农红军第二方面军战史资料选编》（三），科学出版社，1994，第122页。}对苏区普遍的群众教育，国民党方面评判：

\begin{quote}

匪党之所谓“文化教育”、“提高工农群众文化水平”，其麻醉力量较任何宣传煽惑为尤大。盖以邪说灌输脑海之中，改造群众之心理，潜形默化，卒至相率盲从，日趋危途，甘受欺骗而不自觉。尝观匪列宁室墙报处之“识字竞赛”，其进步之程度与麻醉之力量，殊可惊异。\footnote{《新丰特别区政治局局长刘千俊报告匪区民众根本动摇情形匪之维持残局原因及所拟对策》，《军政旬刊》第7期，1933年12月20日。}
\end{quote}


中共在教育上取得的成效，当年多种材料均有反映。据福建上杭才溪区的调查，1934年1月，全区8782人中，除小孩外有6400余人，能看《斗争》的约有8%，能看《红色中华》与写浅白信的约有6%，能看路条与打条子的约有8%，能识50至100字的约占30%，完全不识字的只有10%。\footnote{《才溪消灭文盲运动成绩》，《青年实话》第3卷第8号，1934年2月。}在一个文盲占90%以上的乡村地区，短时期内取得这样的成绩实属不易，任职于国民党中央的王子壮当年曾在日记中写道：“今晨在党部黄君禹同志考察江西匪区归来谈在彼所见云：匪区民众教育的确较普通发达，人民知识程度以高数载。赤匪占据数载而能维持者，悉在组织民众之作用发挥尽致。”\footnote{《王子壮日记》，第2册，台北，中研院近代史研究所，2001，第266页。}


中共对教育的高度重视，既为提高人民的文化水平，丰富人们的精神生活，打破专制政治的愚民统治，也是为了正面宣传中共的政治理念，加深民众对共产主义的理解。因此，中共的教育具有浓厚的政治色彩，采用的方法也是多种多样。红色歌谣就是中共潜移默化影响民众思想的一个重要方法：“仿乡村最通俗的曲调，作成最浅近的歌文，散发各乡村告农民妇女小孩们去唱。”\footnote{《中共江西省委转录赣西各县及二团给赣西特委的报告》，《江西革命历史文件汇集（1929年）》（一），第213页。}“常仿民歌山曲，将其所需要，用以麻醉人民者编为教科书、宣传品，迫令妇女儿童，持诵成熟。”\footnote{赵可师：《赣西收复区各县考察记》，《江西教育旬刊》第10卷第4、5期合刊，1934年7月11日。}如《十骂反革命》、《十骂国民党》以及《十骂蒋介石》等歌曲在苏区被广泛传唱。在此背景下，苏区教育的实用性、宣传性要重于基本知识的输送，教育界主张的“十年树人的计划”就被批评为：“离开了当前的革命需要而从事和平的建设倾向，是目前文化教育工作主要的错误。”\footnote{《中共江西苏区省委四个月（一月至四月）工作总报告》，《江西革命历史文件汇集（1932年）》（一），第159页。}


苏区教育重视现实功用的倾向，和其处身环境的不确定性大有关联，虽然中共成立了苏维埃中央政府，但毕竟宣传意义远大于实际，还不可能与南京政府全面抗衡。在自身无法准确把握未来时，某种程度上的短期行为应可理解。而且，在资源短缺的战争背景下，苏区教师缺乏，经费短少，“每个教师三元伙食费还不可得，而中央政府颁布了对教员的待遇的条例，各县限于经济而不能实施”。\footnote{《中共江西苏区省委四个月（一月至四月）工作总报告》，《江西革命历史文件汇集（1932年）》（一），第160页。}在这样的艰难环境下，苏区教育文化方面的成绩已实属难得，无法过于苛求。


需要提出讨论的仍然是苏区的知识分子政策。发展苏区的文化教育，必须依靠知识分子，但如同中共党内知识分子遭遇洗刷一样，苏区内的知识分子群体处境一度也不容乐观。张国焘在其主持的鄂豫皖中央分局就明确表示：“工农同志在工作中犯了错误党可以原谅三分，倘是知识分子犯了错误，就要加重三分。”\footnote{张国焘：《关于过去的工作和今后的任务（1931年6月30日）》，《鄂豫皖革命根据地》第1册，第305页。}中央苏区虽然话没有说得这么直白，但基本处理方法大同小异。张闻天所说的“吃知识分子”\footnote{洛甫：《论苏维埃政权的文化教育政策》，《斗争》第26期，1933年9月15日。}的现象普遍存在。与此相关，当国民党军进攻苏区时，教育界反水现象相当严重：“广昌县教育部长雷德胜，在敌人还没有来到的时候，就实行叛变逃到白区。接着，城市区的教育部长、甘竹区的教育部长、长生区的正副部长都反水了。而且城市区的教员等有十多人投降敌人”。\footnote{瞿秋白：《阶级战争中的教育》，《斗争》第62期，1934年6月2日。}国民党方面一份关于江西各地在反抗革命中殒命人士材料显示，在所列65人中，乡村知识分子达35人，超过半数；35人中受过小学教育者10人，中学（含师范等中等专业教育）24人，大学预科1人；从职业上看，主要是两大类，一为教育，有12人是县教育局长、中小学校长和教员，一为政治，15人为县区党部人员。\footnote{根据《江西剿匪史料·剿匪阵亡将士及殉难忠烈事略》（1935年编印）统计。}胜利县1932年8月在押犯人621人中，腐朽文人达312人，占总数的50%多。其他包括豪绅地主90人，富农62人，流氓33人，中农33人，贫农88人，工人6人。\footnote{《中共胜利县委八月份工作报告大纲（1932年9月9日）》，《江西革命历史文件汇集（1932年）》（二），第33页。}这些数据显示出知识分子与苏维埃政权间存在相当距离，而这难免会对苏区的文化教育构成消极影响。

\section{社会革命的宣传与实践}


作为一场全方位的革命运动，苏维埃革命在打击和摧毁旧的政治制度同时，对传统社会思想和习俗也在进行着彻底改造。自中共开展革命运动以来，移风易俗就是其革命宣传和实践的重要内容。苏区时期，在群众运动上已经积累丰富经验的中共展开社会革命和群众动员大有驾轻就熟之势。中共的社会革命不是依靠简单的武力压服，更多的是通过简单明了、群众易于接受的宣传动员，改变群众的观念，再以政治力的逐渐渗透，培植新的社会秩序，确立革命的社会基础。


宣传动员是中共社会革命的重要一环。中央苏区的宣传动员形式多样活泼，丰富多彩的群众性文艺活动是中共擅长的方法。群众文艺活动主要有红色歌谣，工农剧社和蓝衫团（苏维埃剧团），以及为群众文化活动提供场所的俱乐部运动。苏维埃剧团各县皆有，而俱乐部也相当普遍。中央苏区规定每乡应设俱乐部，“俱乐部不是少数会弹琴唱曲人的机关，而是广大群众集中的地方，教育群众的集会场所”。\footnote{《闽西各县区文委联席会决议案（1931年7月8日）》，《中央苏区革命文化史料汇编》，第150页。}1934年初，江西、福建、粤赣及瑞金共有俱乐部1917个，参加俱乐部活动的固定会员达9.3万余人。\footnote{《苏区教育的发展》，《红色中华》第239期，1934年9月29日。}传统社会中，民众娱乐活动十分欠缺，看戏几乎是他们最高的精神享受，俱乐部和戏剧演出寓教于乐，可以达到事半功倍的效果。当时，“群众对俱乐部工作表示热烈，一切费用，都是由群众自愿捐助，同志都参加俱乐部各种组织”。\footnote{《一个模范支部的工作报告》，《斗争》第5期，1933年3月15日。}“当看剧团公演时，总是挤得水泄不通。老的小的，男的女的，晚上打着火把，小的替老的搬着凳子，成群结队的来看，最远的有路隔15里或20里的”。\footnote{《苏维埃剧团春耕巡回表演纪事》，《红色中华》第180期，1934年4月26日。}


群众大会是中共开展大规模宣传的另一便捷手段，是各级组织经常性的工作：“每月中至少有四五次大的示威游行，如攻吉、反帝、分土地、成立各级苏维埃、庆祝新年（阴阳历都举行）、欢迎红军、祝捷、慰劳红军、追悼死难者、追悼阵亡战士、成立各军或团、出征等类，不能计其次数。每次各乡区分别举行，人数辄在一、二万，少亦数千”。\footnote{《张怀万巡视赣西南报告》，《中央革命根据地史料选编》（上），第189页。}对于身处僻乡，文化生活极为贫乏的农民而言，能够参与到一些群众性的聚会和游艺活动中，自然十分开心，福建连城“儿童团举行了一次游艺大会，群众散后极高兴，称之为古来未有的事，连旧头脑的人都称为新世界兴旺的征兆”。\footnote{《团闽粤赣省委三个月工作报告》，《闽粤赣革命历史文件汇集（1932～1933）》，第167页。}


苏区报纸、书刊的发行量与群众的文化追求成正比地同步发展。中央苏区从没有报刊到创办有大小报刊34种，《红色中华》从最初的3000份增加到4万份。1932年福建上杭设有农村《青年实话》代售处43所，每期可销售800份。\footnote{《团闽粤赣省委三个月工作报告》，《闽粤赣革命历史文件汇集（1932～1933）》，第169页。}苏区群众体育运动也和文艺运动一样迅速发展，许多地方开辟了运动场，一些“偏（远）乡村中也有了田径赛”。\footnote{《中华苏维埃共和国中央执行委员会与人民委员会对第二次全国苏维埃代表大会的报告》，《中央革命根据地史料选编》（下），第331页。}在战争间隙中，苏区曾举办过全苏区的大规模体育运动大会。


中共通过教育、宣传展开的社会动员成效明显，从当时苏区人民的家书中，可以看出中共革命文化对民众的影响。红军战士史芳森母亲给他的家书中写道：“望你在队，要遵守纪律，用心消灭人敌，我们家中的禾米，有人帮助割回，一切事务，你可不必挂心。”\footnote{《母周氏致儿函》，石叟档案，008·238/2011/0193。}末尾的格式是“此致，赤礼”，完全是标准的同志式语气。曾氏兄弟一为苏区普通民众，一为红军战士，他们的通信中也有如下字句：“家中的生产禾稻，乡府负责人实行派人优待。你不必挂念。现在天气寒冷，你要保障自己的身体，正是革命向前发展，天天进攻敌人。”\footnote{《兄曾衍长致弟曾衍源函》，石叟档案，008·238/2011/0193。}这些信虽然可能是请人代书，但苏区语境的变化还是可以清晰见出。


国民党方面占领苏区后所作调查也证明着中共宣传的力量。莲花县被国民党军占领初期，明显可看到中共宣传动员的影响：“街头巷尾，时见儿童自筑碉堡，隙身其内，而后与碉外儿童互相掷石。若堡外儿童掷石稍稀，碉内儿童即奋勇冲出，与碉外儿童肉搏，非至头破血流不止。更有三三两两，集合一处，讨论对付回籍儿童之方法，其匪化程度之深，于此可见。”“一般人民，无论男女老幼，咸已不复知有民国，年月之计算，仅知公历”。\footnote{赵可师：《赣西收复区各县考察记》（三），《江西教育旬刊》第10卷第4、5期合刊，1934年7月11日。}这种情况在兴国等地也有反映，国民党方面报告，该县“一般男女匪民，只知公历为某某年，而不知民国年号，只知有马克斯、列宁，而不知其它，麻醉之烈，匪化之深，于此可见”。\footnote{《江西兴国县收复后六个月实况报告书》，江西省兴国县档案馆馆藏档案，131/2-8-2/77。}关于苏区革命文化形成的社会思想变化，《申报》记者陈赓雅有如下记载：

\begin{quote}

曾经赤化之人民，似具特性三点：一为不怕死，前几次国军进剿时，彼等皆远飏数十里外，鲜有敢冒险归来者。今则稍逃村外，微知国军能宽容，即联袂偕归。亦有敢越重围寻仇雠，以取甫归之难民首级者。浯塘村民，因割电线，曾被杀数十人以示儆，彼犹不怕，竟乘某连撤退之不备，以农具缴其枪，致遭血洗之祸。……二为残忍性，杀人不算事，其例不胜枚举。三为创造性，譬之义勇队队长，及其它团体主要职位，农民竟敢与农民争夺，争得之后，对厥职亦多能自出心裁，处之裕如；争而失败，则一变而为捣乱心，使办事人时感棘手。\footnote{陈赓雅：《赣皖湘鄂视察记》，第52页。}
\end{quote}


苏区社会革命以扫荡社会陋习为职志，在这方面，禁烟运动的成就相当突出。近代以来，鸦片是危害中国社会的一个毒瘤，南京政府虽然宣称要禁绝鸦片，但实际是禁而不绝，禁烟成为其征税的一种手段。苏区禁烟可谓雷厉风行，对栽种和吸食烟毒者都规定了严厉的惩处措施，并定期展开检查。江西省苏1932年11月报告，在所调查的9个县中，宁都、公略、寻乌无专项调查，其余6县都未发现鸦片栽种情况。兴国、赣县、宁都、永丰、寻乌、安远、万太、公略等地发现少量吸食鸦片者，其中赣县情况比较严重，“吃烟赌博者不少，每月少队所捉的三十多起”。\footnote{《江西省苏报告》，《中央革命根据地史料选编》（下），第236页。}可见，经过中共的努力，在短时期内苏区已经做到基本禁绝鸦片种植，吸食鸦片者大幅减少，残余者也不断采取措施强制戒绝。


由于鸦片吸食后易上瘾，种植鸦片又有经济利益，禁绝鸦片的困难非比寻常，中央苏区为禁绝鸦片，不仅建设并投入警力禁烟抓赌，同时利用少先队等群众组织抓捕瘾君子强制戒烟，经常性的检查也坚持不懈，防止烟毒重新泛滥。1932年底，各地在禁绝鸦片活动中发现胜利、宁都、于都、兴国、石城等县尚有人栽种鸦片，相关各县对此均予以处理：胜利县“进行以政治的鼓动和鸦片恶毒，深入到乡村中去宣传群众、解释群众、发动群众的斗争，把遗下未铲的鸦片，连根铲绝”；\footnote{《中共胜利县委十一月份工作报告大纲（1932年12月10日）》，《江西革命历史文件汇集（1932年）》（二），第379页。}石城县“鸦片烟正在开始铲除”。\footnote{《中共石城县委两个月（10.20～12.20）冲锋工作报告（1932年12月23日）》，《江西革命历史文件汇集（1932年）》（二），第444页。}各级机关派出的巡视员将烟毒列入巡视对象，长生视察于都、胜利两县后就提出报告：“古龙岗、马鞍石一带还有种鸦片的，党与政府亦未注意这一工作”。\footnote{《长生同志关于胜利、于都两县动员工作情况给弼时的报告（1932年12月26日）》，《江西革命历史文件汇集（1932年）》（二），第507页。}通过中共持续不断的努力，虽然在苏区内部烟毒不可能完全禁绝，如当时报告提到：苏维埃干部中吸食鸦片者就有人在，福建团组织发现“上杭、长汀、新泉都有区委负责人偷吃的现象……省委几位十七岁的老人家（原文如此——引者注）也同样吃的”，\footnote{《团闽粤赣省委三个月工作报告》，《闽粤赣革命历史文件汇集（1932～1933）》，第169页。}永新也有“少数党员吸大烟赌钱”，\footnote{《中共湘赣省委关于三个月工作竞赛条约给中央局的总报告》，《湘赣革命根据地》（上），第351页。}但在整个中国烟毒泛滥的背景下，苏区能够在短时期内使种植和吸食鸦片者成为过街老鼠，其禁绝鸦片的决心与社会改造的能力，不能不让人刮目相看。


在艰难的条件下，中共通过社会运动的形式尽力改善苏区的卫生状况。赣南、闽西医疗卫生落后，苏维埃革命前，农民一般病痛多依靠草药治疗。由于医疗条件落后，加之缺乏卫生习惯，常有疫病流行，每年因疫病导致死亡的人口都不在少数：“疾病是苏区中一大仇敌，因为它减弱我们的革命力量。”\footnote{毛泽东：《长冈乡调查》，《毛泽东农村调查文集》，第321页。}苏维埃政权成立后，注意宣传卫生清洁，着力改善公共卫生，兴国县长冈乡将居民按住所接近，七八家左右编为一个卫生班，设班长，负责清洁卫生。规定五天大扫除一次，使卫生状况大有改观。同时政府加强医疗管理，设立公共看病所，“由政府雇请医生，看病不收诊费，但药须在该医生药店购买。规定医生看病费五里以内红包一毛半，每加五里加小洋一毛”。\footnote{《△△区苏维埃政府会议记录》，《江西革命历史文件汇集（补遗部分）》，第482～483页。}1934年春夏，即使在第五次反“围剿”的军事紧张时期，各级政权仍大力开展卫生防疫活动，陈诚在占领石城后发现：“石城一切尚好，较之非匪区尤觉整齐清洁。土匪精神实可令人敬佩。”\footnote{《匪有由信丰安远间突围西窜说深恐一波未平一波又起（1934年10月22日）》，《陈诚先生书信集·家书》（上），第287页。}这应该是由衷之言。


由于时间短暂，加上第五次反“围剿”开始前后苏区社会政治经济状况的恶化，中共中央指导中的教条倾向，苏区社会改革中也有不尽如人意之处。以宣传言，一般来看，像游艺晚会、俱乐部这样寓教于乐的宣传形式往往较能为民众所接受，而一些民众难有切身感受的活动及与群众利益不发生直接关系的组织则无法深入民心。如当时为动员民众，苏区成立了名目繁多的各种群众组织，但效用并不显著：“反帝同盟拥苏大同盟及青年部等革命群众团体的名目太多，每人有十余种组织可加入。下层同志确听到头晕，找头绪找不到。因此这些组织不能健全起来而是一种空招牌。”\footnote{《新泉县委书记杨文仲给福建省委的信》，转见博古《拥护党的布尔雪维克的进攻路线》，《斗争》第3期，1933年2月23日。}1932年底，宁化“全县成立了县工会筹备处、雇农工会、‘反帝’、‘拥苏’、‘互济’等，但组织有的只有一个名，负责人没有，工作更没有”。\footnote{《霍步青给稼蔷同志信（1932年11月8日）》，《闽粤赣革命历史文件汇集（1932～1933）》，第269页。}尤其当缺乏实质内容的会议不断增多时，民众更会产生厌烦心理：“群众大会开多了，群众不愿意到会，且厌恶起来。”\footnote{《张启龙关于赣西工作情形给中央局的报告（1931年7月12日）》，《江西革命历史文件汇集（1931年）》，第107页。}各项工作均走在前列的兴国也报告，在宁都暴动纪念大会中，虽然各级部门大力宣传，但“各乡到会的人数及群众的表现没有很充分的热情”。\footnote{《中共兴国县委两个月冲锋工作报告（1932年12月20日）》，《江西革命历史文件汇集（1932年）》（二），第408页。}这些问题，在国民党方面的报告中也可看到，李一之列举了中共宣传的一些弱点：

\begin{quote}

一、文字宣传，词义艰深，民众无由领会：除标语外，其余各种传单告民众书等，无不千篇一律，长篇大论有如前代八股，宣传效力，直等于零。二、不分时间空间一律施用不能适合环境：如在民族主义高潮中而高呼“武装拥护苏联”，适足引起爱国民众之反感；在农村而高呼“工人实行八小时工作制！”、“工人增加工资”，不啻无的放矢。\footnote{李一之：《剿共随军日记》，第103页。}
\end{quote}


应该说，这确实说到了问题的要害。


苏区对宗教和民间传统信仰习俗采取严厉的态度。《宪法大纲》规定：“一切宗教不能得到苏维埃国家的任何保护和供给费用，一切苏维埃公民有反宗教的宣传的自由。”\footnote{《中华苏维埃共和国宪法大纲》，《苏维埃中国》，第20页。}僧侣在苏区和官僚、地主、豪绅一起，被列入“反革命分子”行列。不过，苏维埃政权打击宗教主要指向的还是组织完整、被认为具有侵略性的西方教会，对中国传统佛、道和民间信仰相对较为宽松。包括和尚、道士以及一些介于宗教和民间信仰活动之间的人士，中共将其定位为游民，属于改造对象，被限制转业，不在被清除的反革命之列。根据毛泽东1933年底的调查，“菩萨庙宇、僧尼道士、算命瞎子等，都成了过去的陈迹，现在是都改业了”。\footnote{《中共赣东北省委代表向中央的报告（1932年11月20日）》，《闽浙赣革命根据地财政经济史料选编》，厦门大学出版社，1988，第131～132页。}


然而，绵延千年的民间信仰活动，短短几年时间想要完全禁绝谈何容易。江西、福建都是民间信仰色彩浓厚地区，江西1930年代的统计，23县中就有寺庙3915所。\footnote{《县政调查统计·江西省》，《内政调查统计表》第22期，1935年6月。}福建40县同期有寺庙2607所。\footnote{《县政调查统计·福建省》，《内政调查统计表》第21期，1935年5月。}赣南、闽西地处山区，民间信仰空气更为浓厚，神佛相当程度上充当着民众应对灾变、祈求保护精神支柱的作用，生命力之强非武力和政治宣传所能及。毛泽东在兴国长冈的调查注意到：“去年以来，老婆太敬神（装香供饭、求神拜佛）的完全没有了，但‘叫魂’的每村还有个把两个……但有些老婆太，虽不敢公开敬神，心里还是信神，这些人多属没有儿子的。”\footnote{毛泽东：《长冈乡调查》，《毛泽东农村调查文集》，第325页。}毛的调查是在工作先进的地区展开的，农村的实际状况可能比这还要严重。湘赣苏区反映：“封建迷信可说大部分打破了，但各县的妇女一遇疾病发生则求神拜佛呼魂的怪象。”\footnote{《湘赣省委妇女部报告（1933年1月7日）》，《江西苏区妇女运动史料选编（1927～1935）》，江西人民出版社，1982，第277页。}莲花县在革命开展数年后发现：“封建迷信的残余，非但不能破除反而逐渐的恢复了……群众对这个问题信得非常浓厚。”\footnote{《少共莲花县委一个月工作计划》，《江西革命历史文件汇集（1932年）》（一），第24页。}


民间信仰旺盛的比较典型案例是瑞金壬田的祭拜观音事件。根据《红色中华》报道，1933年11月22日，壬田区竹塘乡草鞋坪一块白石崩下，岩石崩裂处涌出一股山水，由此被一般信众说成是观音太太，引起乡民祭拜，乃至“我们队伍中一些同志，竟也迷信‘观音’，去烧香顶拜”。\footnote{李忠：《草鞋坪的反革命事件》，《青年实话》第3卷第15期，1934年3月。}短时间内，消息迅速传开，祭拜者络绎不绝，有人还装扮仙童，出卖神签，募捐修庙：“壬田岭子脑地方一个屋子里只有二十多家，一下子就募了四十余块钱。并吸引石城会昌等地的落后群众来朝香，一时轰动远近……许多群众成群结队的前去朝香。”“隘前区苏主席也骑马去拜‘观音太太’，合龙区和云集区的裁判部长也放弃了自己的责任去烧香，还有洋溪乡的乡苏主席，带了二十多个儿童团员前去朝拜”。\footnote{《瑞金壬田区反革命活动》，《红色中华》第160期，1934年3月10日。}壬田离瑞金城不远，堪称中央苏区的中心区，一旦有“灵异”现象出现，朝拜者仍然难以遏制，可见传统习俗信仰在民间扎根之深。在浓重的习俗观念影响下，即使是表现比较先进的青年人也“还有不少信菩萨的”。\footnote{《团瑞金县委关于各区第二次革命竞赛（四月至五月）总结报告》，《江西革命历史文件汇集（1932年）》（一），第294页。}


有意思的是，对苏区传统信仰习俗禁而不绝的局面，国民党人也有评论，杨永泰在国民党军占领广昌后发现：

\begin{quote}

土地庙本是一种迷信的东西，在我们革命势力之下，看见了这个东西，还要劝导一般愚民不可迷信，有些地方党部，激烈一点，便领住些人去打倒他。可是我在广昌黎川经过的地方，看见土地庙土地牌，和香炉存在的就很不少。到处死了人还是化纸钱烧香。听说广昌战后，锡箔纸钱香烛，就是一宗很好的生意，和久缺的盐米一样的需要。黎川还有一座扶乩的菩萨庙，香火从来很盛，偶像陈列得很多，此次经我们收复之后，才把他的偶像搬开去。\footnote{杨永泰：《革命先革心，变政先变俗》，《新生活周刊》第1卷第15期，1934年8月6日。}
\end{quote}


应该说，杨永泰的批评当非杜撰，不过客观来看，这并不能算是忠厚之言。社会改革之功需要一点一滴持续不断地努力，不是单纯政治力量的推动就可以一劳永逸的，对于中国长期的习俗传统，任何毕其功于一役的做法都未免冒进，何况民间信仰还有其更深的社会和心理原因。事实上，后来中共对这一点也有认识，1940年代末，中共在谈到这一问题时就指出：

\begin{quote}

在农民迷信未打破以前，烧神像与佛堂，只有引起农民反感，给反动派以挑拨机会，此外一点好处也没有。过去大革命时期与土地革命时期烧神运动所引起的恶果，今天再不要重复他。今天对农民根本不要提打破迷信问题，提出来农民是不会接受的，何必多此一举？要到什么时候才可以宣传呢？要到农民得到土地，工人生活改善以后，才可以拿事实去教育农民，告诉他们：“菩萨是不会帮助我们得到土地的，什么坏八字，命不好，都是地主欺骗蒙蔽我们的。”只有在这个时候，用这种说法去教育农民，农民才会接受。但就在此时，也不要叫农民不去敬神，只告诉他们不必多花钱就够了。总而言之，这些问题今天对革命尚不是重要问题，不必强调他，自找麻烦。\footnote{《华中分局关于摧毁封建势力及迷信组织给二地委的信（1946年7月7日）》，中共中央华东局《斗争生活》第40期，1946年10月1日。}
\end{quote}


上述总结是中共经过多年摸索后的经验之谈。对于中央苏区时期尚很年轻的中共党人而言，他们考虑更多的还是革命理想的实践，社会革命是共产革命的题中应有之义，而中共党人通过社会革命在实现自己革命理想的同时，也相当程度上赢得了民众的好感与支持。然而，正如中共自己后来意识到的，由于社会的运行呈现出无限多样的复杂性，社会革命的掌控和把握相当微妙，其效果往往利弊参半。比如，苏区宣传在家庭中打破家长制度，家庭成员平等，就激起一些中老年人的反激，苏区民众不同年龄段对中共支持度出现波动，与此不无关系。资料显示，中央苏区老、中、青、幼各年龄段对中共反应有相当差别。福建方面有报告提到：“思想方面，则老年人痛恨赤匪，冀得真命天子出而恢复专制，平治天下。幼稚者因受共匪愚惑，以为非阶级斗争，实行土地革命，决无其它出路。惟少数中年人稍能折中于两者之间，希望有不杀人之政府出现，使彼有安居乐业之机会，即为己足。”\footnote{李大奎：《崇安县教育调查报告》，《福建教育周刊》第213期，1934年11月26日。}这种随着年龄增长对苏维埃革命支持度呈递减趋势的状况在江西方面也得到证实：

\begin{quote}

（老年人）四十岁以上的属之，他们阅历较深，态度稳健，除了一些素行不轨的以外，大都不赞成土匪的行动，在土匪一方面，也怀疑他们思想顽固，目为反动。所以在匪区被杀的，亦以老年而有知识的人为最多。但无论如何，他们在群众中，对于后进，有相当信仰，在每一时期，因为环境的压迫，不能不和土匪虚与委蛇，然而他们处在岌岌不可终日的境况中，心理上的感应是惊恐、怨望、饮恨，绝少对土匪以助力。

（壮年人）二十岁以上四十岁以下的属之，这一班人，血气方刚，意志未定，一经刺激，不问事之如何，辄生反应，故麻醉较易。且好新奇，喜破坏，土匪之所以蓬勃一时，实以利用此辈之力为多，这可以说是土匪的主力军。但据各次战役所俘获的壮年匪兵所供称，他们之所以不能投诚，完全是由于土匪政委秘密的监视。可知他们的意识，也发生了动摇。并且一次二次的抽派到前线抵炮火，所有的壮丁，已十去八九了。

（幼年人）十岁以上廿岁以下的属之，这一期内的人，口臭未退，初出茅庐，社会上人情世故，根本不懂，而破坏好奇的性格，又较壮年为甚，所以极易利用。最近在樟村之役，俘获匪兵，多系十余岁的儿童，由匪区所抽派者，他们除了感觉痛苦以外，什么也不明白，有无数儿童，就是这样挡着炮火糊糊涂涂死了，所以在幼年儿童方面，他们根本没有成见，也无所谓觉悟。\footnote{欧阳明：《匪区民众心理之分析与挽救》，《汗血月刊》第1卷第2、3期合刊，1934年2月20日。}
\end{quote}


年龄差异导致的支持度变化，除观念的原因外，和其在社会革命中面临的地位变化当然不会没有关系，老年人革命后从家长的权威中跌落，而青少年的束缚相应得到减轻，这是导致他们对革命后的社会变化反应不同的关键原因，这一点，妇女群体体现得尤为明显。

\section{妇女：地位上升最快群体}


马克思主义的共产革命以求平等、求解放为宗旨，具有强烈的弱势代言人色彩，而妇女无论中外，在社会生活中都长期居于弱势，是革命的天然拯救对象。恩格斯引述傅立叶时曾说道：“在任何社会中，妇女的解放程度是衡量普遍解放的天然尺度。”\footnote{恩格斯：《反杜林论》，《马克思恩格斯选集》第3卷，人民出版社，1995，第610页。}妇女于社会革命中的这种独特地位，在苏维埃革命中也有清晰体现，苏维埃时期所有的社会变化中，妇女地位的改变可以说最为引人注目。


赣南、闽西是客家人聚居地，客家人有天足传统，在一些较为偏僻地区，客家妇女天足比例很高。由于不受小脚限制，加之客家文化相对于中国主流文化的边缘地位，客家妇女历来在家庭劳动中占有重要地位，尤其在离城较远的偏远地区妇女参加劳动比例更高。靠近赣南的井冈山一带客家人比例较高地区，妇女参加劳动者占相当比重，资料记载，湘赣苏区的劳动妇女，“除（北路攸县萍乡）一部分是小足外，其余各县的劳动妇女都能参加农村生产做事耕田”。\footnote{《湘赣省委妇女部报告（1933年1月7日）》，《江西苏区妇女运动史料选编（1927～1935）》，第276页。}莲花“妇女向系天足，其操作与男子等（缠足者极少）”。\footnote{赵可师：《赣西收复区各县考察记》（三），《江西教育旬刊》第10卷第4、5期合刊，1934年7月11日。}


赣南、闽西客家妇女参加劳动者甚众。时人报告：“闽西各县除龙岩外，妇女向来天足，能自食其力（如斫柴田事多妇女担任），农人工人对于家庭负担自属轻微，且多妇女能苦力生产维持一家而男子毫不负责者。”\footnote{《漳州永定救乡委员熊耀球呈国民政府（1930年11月16日）》，《伪国民政府部分档案中有关第二次国内革命战争史料抄件》第6册，中国社会科学院近代史研究所藏，第571页。}关于赣南妇女天足的报告更多，“在赣县女人是生产的主力，这并不是怎样夸大的话。乡间种田的是女人，码头上以及长途挑运的人也是女人”；\footnote{沈锡忠：《我们知道的赣县和赣县人》，《新赣南旬刊》第1期，1941年2月1日。}南康“妇女尽属天足，劳动工作，不逊男子，且均能自食其力，绝少为人奴仆者”。\footnote{《南康县教育视察报告》，《江西教育》第10期，1935年8月1日。}客家妇女以参加辅助劳动为多，但也有部分妇女进行田间耕作，革命前上才溪妇女中“约三十个人能用牛”。\footnote{毛泽东：《才溪乡调查》，《毛泽东农村调查文集》，第342页。}


妇女勤于劳作，相当程度上减轻了男人的负担，妇女可以不依赖男人，甚至养活男人：“赣南妇女都是和男人一样的大脚，耕田做工都是和男人一样的负担，一般的以为一个女子能自己做事养活男子是光荣。”\footnote{《黎日晖关于赣南工作的综合报告（1931年10月6日）》，《江西革命历史文件汇集（1931年）》，第183页。}由于妇女的辛劳操作，在赣南的上犹县甚至出现阴盛阳衰的状况：

\begin{quote}

十二万多人口的上犹，女性是占了七成多。所以各项工作都有女性做。尤其是劳动方面，妇女是其主力……但是上犹的男子是太闲了，太懦弱了。像这一个月里的天气，实在是并不很冷的，可是穿着长衫，在长衫下提个火笼到街上走的男子，是常看见的事。\footnote{魏晋：《建设中的上犹》，《新赣南旬刊》第5卷第3期，1943年2月15日。}
\end{quote}


当然，上述材料提供的妇女参加劳动较多状况应该是相对其他地区而言，尽管客家妇女有天足传统，但流风所及，赣南闽西妇女缠足风气仍然存在。中共方面文件显示：“赣东（如宁都、石城、广昌）等县……女子在十六岁以上的，大部分还是小脚，劳动力弱。”\footnote{《中共江西苏区省委四个月（一月至四月）工作总报告（1932年5月）》，《江西革命历史文件汇集（1932年）》（一），第168页。}胜利县的“古龙岗、江口、赖村、半迳等处的小足妇女也实在多……大概还有十分之八小足婆”。\footnote{《中共胜利县委关于七月份工作情形给中共中央局的报告》，《江西革命历史文件汇集（1932年）》（一），第379页。}一般而言，客家人口比例愈高，地区愈偏僻，则天足比例愈高，而在城镇及通衢大道上，因为受流行习俗压力较大，缠足比例相对较高。\footnote{考量这一点，必须注意到1930年代前后放足运动开展仅十多年，成年妇女的缠足状况和放足运动尚难形成相关，因此在赣南、闽西客家区域的城镇地区小足反而较多。}


中央苏区妇女特殊的天足状况，一方面使妇女与外界的联系更为密切，在家中承担更多责任，由此带来一定的社会和经济地位；但同时也使这里的妇女须承受更多的生活压力：“赣西南的妇女有百分之九十以上是农村中的劳动妇女。她的生活、工作，一切都比男人特别利害，尤其是青年妇女更受痛苦。因为在那边古来的习惯，一般女子除从男人出去耕田外，她回家要做饭，及家里其它一切事情。”\footnote{《朱昌谐同志关于赣西南妇女工作的报告（1930年10月23日）》，《江西革命历史文件汇集（1930年）》（二），中央档案馆、江西省档案馆编印，1988，第114页。}而且，从总体上看，在当时的社会和文化氛围中，妇女的地位仍然低于男子，“男子压迫妇女，轻女重男的现象……丈夫打骂老婆，家婆打骂媳妇，还是认为天经地义”。\footnote{《江西苏区中共省委工作总结报告》，《中央革命根据地史料选编》（上），第473页。}女性的社会角色使她们自觉不自觉地成为困苦生活的最大牺牲者：“她的生活更比男人要苦得多，比如男人做事回来或者要买点好菜吃的时候（那边大部分男人要吃酒），女子是没有吃的，并且女子吃饭除五十岁以上的老婆婆外，青年女子及成年女子，统统没有资格上台吃饭的，其余穿衣方面，什么一切都比较男人要苦点。”\footnote{《朱昌谐同志关于赣西南妇女工作的报告（1930年10月23日）》，《江西革命历史文件汇集（1930年）》（二），第114页。}中国传统社会习见的妇女社会地位低下的状况并没有根本改变。


1920年代后，中国社会掀起男女平权运动，南京国民政府理论上也提倡男女平等。但是，外界的影响和政府的权威在这里若有若无，妇女地位看不到变动的迹象。中共进入后，秉持男女平等、提高妇女地位的理念，通过强有力的政治推动，大力提高妇女权益。苏区政府颁布法令，禁止虐待童养媳并废除童养媳制度，禁止翁姑丈夫虐待妻媳，禁止幼女缠足、穿环，30岁以下妇女应剪发、放足；严禁贩卖女子，违者枪决；取消蓄婢制度，一切婢女即行解放并由主人给以工资维持生活；女子被压迫为娼妓者，即行解放，恢复一切自由。同时，苏维埃采取男女平等的社会经济政策，所有城乡女工都得到与男工同样的一份劳动工资。农村中妇女、小孩都得到按人口平均分配的一份土地和山林。对分到的土地，“她们均有权自由处置，不受任何人的约束”。\footnote{《列宁青年》第5卷第1期，1931年12月8日。}1932年临时中央政府发布训令，决定从中央到地方建立妇女生活改善委员会，承担调查、统计和研究苏区妇女生活状况，拟定改善妇女生活的建议和办法，考察有关妇女政策、法令执行情况等任务。一系列的法规规定了妇女独立的经济地位，为妇女争得与男子平等的社会经济地位奠定了政策和物质基础。


在传统中国，妇女很少有受教育的机会，赣南、闽西妇女更是如此。教育水平的低下使妇女无法在知识、视野与社会活动能力上与男子齐头并进。苏维埃为妇女提供了享受文化教育的权利和机会。苏维埃中央政府专门要求各级政府文化部设立妇女半日学校，组织妇女识字班、家庭临时训练班、田间流动识字班，有计划地实施妇女的文化工作。规定7至18岁的女子应接受小学教育，19至30岁的入夜校学习，失业女工到妇女工读社受教育，其他妇女也要参加识字运动。通过政府的大力提倡、引导，女子接受教育比率有了大幅提高，兴国县列宁小学的学生中，女生占45%；夜校成人学生中，女子占69%；识字组男女组员中，女子占60%。


婚姻自由是妇女解放的重要内容。传统中国婚姻由家庭包办，男子可以纳妾，妇女在婚姻中处于绝对的弱势地位。苏区颁布的婚姻条例宣布，废除一切封建的包办、强迫与买卖的婚姻制度，实行婚姻自由和一夫一妻制。在离婚问题上，实行“偏于保护女子，而把因离婚而起的义务和责任，多交给男子担负”\footnote{《中华苏维埃共和国婚姻条例》，《中央革命根据地史料选编》（下），第194页。}的优待妇女的政策。新的婚姻政策造就了不一样的婚姻关系，深有所感的赣南民众在山歌中表达出他们的感触：“实实在在话你知，共产主义矛（冇——引者注）共妻，总爱两人心甘愿，唔使媒人也可以。”\footnote{《共产主义矛共妻》，《赣南苏区民歌》，赣南人民出版社，1959，第8页。}


苏维埃政权着力提高妇女地位的措施，收到明显成效。毛泽东在第二次全国苏维埃代表大会上谈道：“兴国等地妇女从文盲中得到了初步的解放，因此妇女的活动十分积极起来。妇女不但自己受教育，而且已在主持教育，许多妇女是在作小学与夜校的校长，作教育委员与识字委员会的委员了。”\footnote{毛泽东：《对第二次全国苏维埃代表大会工作报告》，《苏维埃中国》，第283～284页。}这种地位的变化，在日常生活中处处可见，毛泽东在兴国长冈乡调查发现：“丈夫骂老婆的少，老婆骂丈夫的反倒多起来了。”\footnote{毛泽东：《长冈乡调查》，《毛泽东农村调查文集》，第324页。}


苏区妇女地位的提高，激发了她们参加社会政治经济生活的热情，无论在后方生产还是在安全警戒上都发挥着越来越重要的作用。各地在乡苏维埃政府之下，设立妇女劳动教育委员会，积极推动妇女参加生产劳动，许多妇女掌握了犁耙、莳田技术，担负起原来由男人承担的繁重劳动。1934年春耕期间，瑞金能够参加生产的妇女达到3104人，仅下洲区就有1019人。\footnote{定一：《春耕运动在瑞京》，《斗争》第54期，1934年4月7日。}兴国全县参加生产的妇女更高达两万人以上。\footnote{王首道：《模范红军家属运动》，《斗争》第70期，1934年8月16日。}青年妇女则参加到少先队等群众组织中，成为保卫地方安全不可或缺的力量：

\begin{quote}

上杭大拨区大拨乡的一个妇女细蓝同志，她是一个团员，当敌人前次进攻该乡时，男同志及部分的群众都离开本乡了，但是这位模范的女性却安静沉着的不跑，当敌人快要迫近大拨时，她一面燃烧号炮——因为放哨的人都“逃之夭夭”了，俾广大工农群众知悉准备与敌作战，一面亲自率领模范少队数十人奋勇直前的打击敌人，给了进攻的敌人以严厉的痛击。\footnote{飞涛：《三个模范的女性》，《青年实话》第2卷第10期，1933年4月1日。}
\end{quote}


妇女广泛参加政治生活是妇女地位变化的重要标志之一。当时，“苏维埃政府有女子当选，一切群众示威游行等运动，均有女人参加，作战时妇女送饭茶慰问伤兵都极热烈。新年耍龙灯，女子都提灯、化装、武装出来了”。\footnote{《张怀万巡视赣西南报告》，《中央革命根据地史料选编》（上），第192页。}1933年5月，中共闽赣省委在召集省党代表大会的有关通知中，明确规定：“出席各级代表大会的代表，工人成份应占百分之四十，劳动妇女至少要占百分之十。”\footnote{《中共闽赣省委关于准备召开全省党代表大会的决定》，《闽粤赣革命历史文件汇集（1933～1936年）》，中央档案馆编印，1984，第4页。}在1933年下半年开展的苏维埃代表选举中，苏区妇女积极投身其中，瑞金下肖区记载这里的妇女选举情况：

\begin{quote}

区妇女部接到中央局通知，立即召集妇女委员会，具体的讨论。党团区委常委会通知各支部召集各妇女指导员和主席团联席会议，此会议上党团区委首先具体的讨论选举运动，过后拿到此会中去进行，很迅速地在一个多钟头的会议中，先半个钟头政治与前方红军伟大胜利好消息报告，她们非常兴奋，根据此报告来讨论了下列问题：1.妇女在苏维埃选举百分之二十五的准备。2.给全苏大会赠品，各乡热烈建议扩大红军竞赛，全区共四十多名。3.每个妇女做一只草鞋给前方红军。

对于过去婚姻条例的意见，做提案给全苏大会修改，各乡妇女具体讨论过（如妇女自由带土地财产与小孩子给养费等问题），这任务完成了。由区一级组织了总结各乡婚姻条例提案委员会，全苏大会做修改婚姻条例参考。\footnote{汉英：《瑞京下肖区妇女参加选举运动的检查》，《青年实话》第3卷第1期，1933年11月13日。}
\end{quote}


随着妇女在政治生活中的日渐活跃，妇女广泛参加到各级政权中。1933年，江西苏区16个县，县级妇女干部有27人，兴国一县有20多名妇女担任乡主席。1934年初新组成的城乡苏维埃代表大会中，妇女代表一般占代表总数的25%以上，有的地方如上杭县的上才溪乡和下才溪乡，妇女代表分别占54.6%和64.8%。\footnote{梁柏台：《今年选举的初步总结》，《红色中华》第139期，1934年1月1日。}在各级政权机关中，“乡政府及区、县政府，亦统统有女子当选委员”。\footnote{《朱昌谐关于赣西南妇运的报告（1930年10月23日）》，《赣南妇女运动史料选编》第1册，江西省妇联赣州地区办事处编印，1997，第15页。}1933年底江西省苏维埃政府执行委员候选名单中，妇女人数多达31人。中共党员中妇女党员人数随着妇女运动的活跃日益增长，据统计，江西省女党员在党员总数中的比例，1932年夏为8%，次年1月，据瑞金、宁都、兴国等九县的统计，增加到13%，总数达5871名。\footnote{弼时：《为扩大一倍新党员而斗争》，《斗争》第11期，1933年5月10日。}5月，总数进一步增加到10294名。\footnote{《党的组织状况》，《中央革命根据地史料选编》（上），第674～676页。}对于苏区妇女的政治热情和地位的提高，国民党方面也有客观描述，承认苏区“一般妇女须随时常为红军缝做草鞋，而妇女结婚、离婚，绝对自由”。\footnote{《南昌行营致汪精卫等电（1934年4月13日）》，《中央革命根据地革命与反革命斗争史料》（下），中国社会科学院近代史研究所藏。}


中央苏区妇女的天足及与外界联系较多的状况为中共在此展开妇女运动提供了比较好的天然条件。苏区妇女本身就有着当家养家的传统，其政治热情和参与意识较易得到激发，因此苏区妇女在扩红、慰劳红军乃至后方生产等方面发挥着十分重要的作用。福建长汀扩红运动中，“在群众会上妇女突击队，指男人的名字，唱山歌鼓动他当红军，鼓动逃避归队。在这样有力的突击中，很多男子马上报名。妇女突击队大多是红军的老婆组织起来的，他在群众中，有他极光荣的威信和作用”。\footnote{郭滴人：《长汀最近扩大红军所得的经验》，《闽粤赣革命历史文件汇集（1932～1933）》，第289页。}兴国高兴区妇女“扩大红军占全区半数以上，推销公债占百分之七十”。\footnote{江西十县参观团：《参观兴国以后的感想》，《红的江西》第18期，1933年1月29日。}中共在妇女工作上的成功表现，使国民党军为之头痛，特地指示部队官兵：“匪设有妇女训练班，经过训练后，装难民逃到我军地区求夫，以作兵运，及刺探我方消息，并作种种活动。请转令各部队，不论官兵，不得在匪区结婚，并禁止士兵与妇女接谈。”\footnote{《训令东西北路总司令部令各部队不得在匪区结婚并禁止士兵与妇女接谈》，《军政旬刊》第30、31期合刊，1934年8月20日。}


在充分肯定妇女解放的效果和作用时，也应看到，妇女解放对社会牵涉极大，其过程相当复杂。虽然赣南、闽西的客家妇女和中国主流文化中的妇女有所差别，但男尊女卑仍然是这里的基本构架，妇女挑战千百年来的习俗，抛头露面参与到社会政治活动中，遇到的阻力相当巨大。当时有报告反映：“永定有妇女的组织，一般妇女出来都是很腐化的，没有出来则一点与男子说话不敢，甚至开会怕人家骂的。”\footnote{《闽西同志口头报告（1930年12月19日）》，《福建革命历史文件汇集》甲16册，第113页。}1932年，“妇女工作较好”的兴国县女党员仍只有289人，占全体党员数的7%；宁都1500名党员中只有7名女党员。\footnote{《江西苏区中共省委工作总结报告》，《中央革命根据地史料选编》（上），第491页。}1933年江西报告：“妇女干部数目字的少达到惊人的程度，只占总数百分之六点四。”\footnote{《党的组织状况》，《中央革命根据地史料选编》（上），第687页。}妇女在社会政治中的作用仍远未充分展现。


妇女解放现实阻力来自于异性。对于长期居于权力、家庭主导地位的男子而言，妇女对社会政治生活的强势介入引起他们的抵触实属意料之中：

\begin{quote}

当妇女们自己组织起来，参加集会，参加社会活动时，却越来越遭到男人们，特别是自己的男人的反对。男人大都认为，娘们出去参加活动，一定会引起“伤风败俗的事”来。男人花了不少粮食才娶下媳妇，因此便把她们看做是自己的私产，巴望她们出力干活，生儿育女，伺候自己，伺候公婆，只有别人问话时才许搭腔。在这种气氛下，妇女会的活动使许多家庭都产生了家庭危机。不仅做丈夫的反对自己的女人出门，公婆反对得更厉害。许多年轻媳妇因为坚持出去开会，回家后便遭到毒打。\footnote{〔美〕韩丁：《翻身——中国一个村庄的革命纪实》，第179页。}
\end{quote}


不仅仅是普通农民，即使是一些苏维埃干部，也不能摆脱长期观念熏陶的影响，一些地方政府有意无意中成为抵制妇女解放的堡垒：“永丰有很多地方对婚姻条例完全忽视，有妇女坚决要求离婚，政府不准许外，还要勒逼女子要离婚就要出洋几十元”。\footnote{《江西省苏报告》，《中央革命根据地史料选编》（下），第237页。}“寻乌吉潭区有个乡苏主席，因为同姓的一个女子同人恋爱，要区政府严办那个人而区政府不同意，该乡苏主席竟带领一班游击队，子弹土筒，并发动当地群众捣乱区政府，而县苏裁判部仅判处该乡苏主席一年苦工。”\footnote{司法人民委员部：《对裁判机关工作的指示（1933年6月）》，石叟档案008·548/3449/0759，中国社会科学院近代史研究所藏。}


应该指出，中共妇女政策对于数千年社会结构的冲击，其所造成的社会动荡是巨大的，引起反激势在必然，而初期妇女运动的错误又加剧了反对的力量，并引发一些社会问题。初期妇女运动单纯以婚姻自由相号召，发生不少荒唐现象，各地都有报告：“以‘离婚结婚绝对自由’、‘打破封建观念’等狭隘的整个斗争脱离的口号来动员青年妇女，实际的进行浪漫行动故意的捣乱，如妇女要能当众脱裤才算是真正打破封建观念，如组织恋爱研究会等等”。\footnote{《江西苏区中共省委工作总结报告》，《中央革命根据地史料选编》（上），第472页。}一些干部受所谓“杯水主义”影响，对妇女不负责任：“今天和这个女人结婚，明天又和那个女人结婚，上级机关的负责同志如此，影响到下级群众更加是一塌糊涂”。\footnote{《中共闽粤赣特委报告第三号——政治形势和党的工作（1931年4月7日）》，《闽粤赣革命历史文件汇集（1930～1931）》，第114～115页。}有些妇女也把婚姻解放误为性混乱，造成男女关系混乱，有报告反映：“在安福近来每个女子，特别是所谓开通的都有三个男子，一个是丈夫两个是候补，两个中有一个半公开，一个秘密的，将来，现在的丈夫合不来离开了婚即以公开的递补，另一个秘密的。在西区每个男同志有一个老婆一个爱人一个Ⅹ老婆共三个等的古奇事。”\footnote{《张启龙关于赣西工作情形给中央局的报告（1931年7月12日）》，《江西革命历史文件汇集（1931年）》，第106页。}


福建杭武误解打破封建口号，青年活动时关系混乱：

\begin{quote}

男男女女，扳头拉颈，会后，即男找女，女找男，三个五个，男男女女共睡一床。少先队下操做蛇脱壳，脱裤子，接塔等……假使上述事情谁怕做，谁不愿做，谁就是“封建”，就要受处罚，甚至开除队籍。因为这样来“打破封建”，使得一般青年妇女怕来下操开会，有些群众反对下操开会，以至反对“反对封建”、“男女平等”，对革命不满。\footnote{清洪：《这样打破封建要得吗？》，《青年实话》第12号，1932年3月15日。本文作者清洪即赖清鸿，为少共福建省委巡视员。文章发表后，少共上杭县委致函《青年实话》，表示：“男女杂卧一事，在会议时，有时是有的。譬如赖清鸿同志于二月五日到我区来巡视召集团员大会，于八日开幕，亦因天气太冷，有许多同志没有床睡，杂卧一起，清鸿同志都与女同志共卧了！他并要交头共卧，至半夜连灯火都吹灭了。”见《少共上杭县委来信》，《青年实话》第15号，1932年4月15日。}
\end{quote}


片面提倡妇女解放导致的性错乱现象直接威胁到农民家庭的稳定。在性解放的背景下，普通农民处于最弱势的地位，因为“工作人员凭借政治的地位与权力，的确有假自由恋爱之名，行夺人妻女之实，自然其中有不少女子有趋炎附势的事实。农民则屈于淫威之下，敢怒而不敢言。在则女子也有浪漫放纵生活奢侈者”。\footnote{《邓中夏同志关于红二六军的报告》，《湘鄂西苏区革命历史文件汇集》第2册，第39页。}因此，“一般农民恐惧其已有的老婆被小白脸的知识分子夺去，未婚妻子不肯来到自己家里，将有没老婆睡觉的危险”，\footnote{《张怀万巡视赣西南报告》，《中央革命根据地史料选编》（上），第193页。}对性解放十分反感。即使是中共党员，对妇女出来工作也有保留：“很多共产党员不愿意他的老婆入党，我们推测这种原因，大概是恐怕自己的老婆入党后，又要给人家自由去。”\footnote{《中共闽粤赣特委报告第三号——政治形势和党的工作（1931年4月7日）》，《闽粤赣革命历史文件汇集（1930～1931）》，第114～115页。}由于婚姻危机成为社会群体的恐惧，地方政府对婚姻和男女关系问题稍一处理不当，极易惹起众怒，赣东北横峰县一女团员要与丈夫离婚，因政府处理不当，“村中有三百多农民自动召集大会，全体向苏维埃政府交涉，说女团员借故离婚，男子并未压迫，请苏政府释放，否则，我们村中五百多人都不要老婆了，我们将五十岁以下的老婆不要”。对此，苏维埃政府不得不表示退却，“经苏维埃政府召集会议，党召集党的会议，妇女工作人员会议，决定释放这一男子，处罚女团员”。\footnote{《赣东北妇女工作》，《江西苏区妇女运动史料选编（1927～1935）》，第419页。}


面对上述种种难题，妇女运动经历初期理想主义的浪漫后，终究要回归柴米油盐的现实。1932年，福建永定县委书记萧向荣在《红色中华》发表公开信，就苏区婚姻条例中关于离婚的规定提出几点质疑：假使一个男子或女子，没有一点正当理由提出了离婚，究竟可否准其离婚，“离婚绝对自由”引发的朝秦暮楚现象如何解决；“男女同居所负的公共债务，归男子负责清偿”，假使女子因负债太多要求离婚，男子负担是否过重；“离婚后，女子如未再行结婚，男子须继续其生活，或代耕田地，直至再行结婚为止”，如果女子无理由要求离婚，而离婚后男子还要负担女子的生活费，岂非雪上加霜。\footnote{《问题与答辩》，《红色中华》第11期，1932年2月24日。}当时，苏维埃临时中央政府副主席项英在答复中未能直面这些现实问题，继续空洞地强调保证妇女的离婚自由，似乎，项英没有意识到或者不愿正面面对萧向荣所触及的革命理想和现实操作之间的落差，即尽管婚姻自由是中共努力争取的目标，但矫枉过正的规定又在损害着其他群体的利益，也未必就一定会有好的结果。社会是一个可以自我循环的有机体，当中共开始在一定程度上以执政而不完全是革命的姿态进入社会管理时，将会发现，利益的调节、资源的分配、习俗的影响，这些琐碎而又现实的问题虽然会让理想从云端回到地面，却又不能不认真面对。革命可以荡污涤垢，但社会秩序的重建还是需要一点一滴的改造之功，而且，这中间很可能会经历曲折反复的过程。


所以，现实的状况是，随着初期妇女解放的宣传逐渐转向妇女权益保护、妇女地位提高等具体问题，妇女运动的调门和受重视程度明显降低。在传统观念依然浓厚的背景下，妇女作为一个弱势群体受到的特殊关注一旦削弱，有些初期受到控制、纠正的现象又有重起之势。在男女平等问题上，不少地区反映：“有些乡村封建压迫存在着，如男人打骂妇女（如沙洲乡一个女子，七保乡一个女子），在解决婚姻问题上发生许多严重的事情，当地党和团都往往不去解决。”\footnote{汉英：《瑞京下肖区妇女参加选举运动的检查》，《青年实话》第3卷第1期，1933年11月13日。}婚姻自由虽然在法律上得到了充分的保障，但由于习惯礼俗的影响，也未能完全禁绝，胜利县委报告：“打骂妇女阻止妇女开会做工作和虐待童养媳的，怪事也仍然有的。”\footnote{《胜利县委工作报告（1932年10月8日）》，《赣南妇女运动史料选编》第1册，第57页。}永新县的状况更为严重：

\begin{quote}

苏区的劳动妇女还有不少的受残酷凌辱，至于打骂逼死妇女成了普遍的现象。如南阳区某乡用沸水泡死童养媳，象形区打出童养媳几个月不去寻问，并花溪乡有个童养媳在此严寒酷冷的天气中盖蓑衣……甚至有些妇女在婚姻问题上不能得到婚姻自由，反受到父母及旁人压迫干涉的手段而自寻短见……潞江区厚田乡有个青年妇女为要结婚而父母不准服药而死。\footnote{《永新县苏维埃执行委员会训令（1933年1月31日）》，《江西苏区妇女运动史料选编（1927～1935）》，第289页。}
\end{quote}


诸如童养媳之类的习俗在民众中事实上更多受到习惯法的保护，“如叶坪就有二、三个童养媳，不愿在十五、六岁时，同他老公结婚，更不愿受家婆的压迫和打骂向政府报告，当着政府机关的人去调查时，他们的邻居都以‘女大当嫁、家婆对他满好’来搪塞”。\footnote{伯钊：《纪念“三八”与妇女工作应有的转变》，《红色中华》第12期，1932年3月2日。}其实，这并不一定完全是搪塞，因为在民众心目中，这些确实已经司空见惯，习以为常。甚至有地方政府公然强迫婚姻的：“乡政府或区政府有可以指定女子与某人结婚，男子有时贿通乡政府来达到与某女子结婚的目的。”\footnote{《江西苏区中共省委工作总结报告》，《中央革命根据地史料选编》（上），第473页。}


对长期以来形成的风俗习惯，中共在妇女运动中虽有注意，但离根除尚很远。国民党军占领广昌后发现，“此地因文化落后，教育极不发达之故，其风俗习惯，自是守旧，毫无革新改进可言！故无论城乡妇女，一律大衣小袴，缠足结髻，银簪束起，银牌系后”。\footnote{《广昌社会调查》，《汗血月刊》第1卷第5期，1934年7月20日。}1934年10月的报告则称，宜黄“妇女缠足之风犹盛”。\footnote{《宜黄第七区民训会工作》，《军政旬刊》第37、38期合刊，1934年10月31日。1935年国民党军占领宜黄后，有调查者也发现：“本县少女一至十五六岁，即挽髻缠足，殊有改良之必要耳。”见周炳文《调查日记》，萧铮主编《民国二十年代中国大陆土地问题资料》第172种，第85728页。}对于广大妇女和苏维埃政权而言，中央苏区在男女平等上取得的成就固然喜人，但要完全实现男女平权，还不是一蹴而就那么简单。

\section{群众：组织与改造}


苏维埃革命开展后，中共犹如横空出世，在结构松散的中国农村，建立起一个紧密的、具有有效组织及动员能力的社会体系，令其对手方不能不为之折服。抗战开始后，谢觉哉以中共驻兰州办事处负责人身份在兰州活动，与国民党人接触时，听到国民党人的感叹：“到警署晤马志超局长。马参加过剿共，很惊异苏区群众组织，他们派的侦探，不能进入苏区五里路。”\footnote{《谢觉哉日记》，1937年7月30日，人民出版社，1984，第122页。}


国民党人说到的中共的民众组织，的确是中共组织和动员能力的最好见证。当时有关调查提供了苏区组织的一般状况，工作偏于中上的兴国高兴区调查数据是：“高兴区是兴国县第二等先进区域，虽然比上社区、城市区较差些，然而一般的还算好。这一区有两万零七百多人口，平均男有一万多人，女则不到一万人，党员1026名，团员989名，赤卫军全区有十四连（每连120人），共一千七百多名。工人、雇农不少。少先队有1300人，内有妇女九百多人。儿童团员有2748人。”\footnote{如心：《兴国高兴区两周印象记》，《红校斗争》第5期，1933年7月8日。}数据显示，党员和团员分别占到总人口的近5%，参加赤卫军、少先队、儿童团的民众占到1/4强，这是一个以党员为基干，各种民众组织为补充，层层相连、递相推进的组织体系。


反映着中共广泛动员民众的目标，中央苏区民众组织众多，包括工会、贫农团、妇女会、反帝大同盟等，其中，群众性的武装组织最受重视，这是由苏区处身战争环境所决定的。早在1930年5月，毛泽东以红四军前委名义提出组建农民武装问题，要求在新开辟苏区“数天之内分完田地，组织苏维埃，建立起‘赤卫队网’（所有十六岁以上四十五岁以下的男女壮丁一概编入，由乡到特区各级指挥机关，隶属于各级苏维埃的军事部）……赤卫队不但可以代替农会的作用（团结群众），并且加了一层‘武装起来’的意义”。\footnote{《对流氓和对农民武装的策略——前委给安于会赣四县边界特区的信（1930年5月）》，江西省博物馆存。转见何友良《苏区制度、社会与民众研究》，未刊稿，第121页。}1932年7月，中华苏维埃共和国中央执行委员会向中央苏区各县发布训令，指示苏区“各县或几县成立一个红军补充团”，同时“立即普遍的发展和成立城乡赤卫军”，\footnote{《中华苏维埃共和国中央执行委员会训令（第14号）》，《中央革命根据地史料选编》（中），第620～622页。}以此作为广大群众的武装组织和红军的补充队、后备军。1932年9月，中华苏维埃临时中央政府又要求各县“赤卫军每区成立模范营，每县成立一个模范师，以统一指挥。不仅是随时集中配合红军行动，巩固和发展苏区，并作赤卫军下级的干部训练地方，和准备到补充队和直接到红军的积极分子组织”。\footnote{《中央执行委员会关于扩大红军问题训令（第15号）》，《中央革命根据地史料选编》（中），第641页。}


苏区群众武装组织体系严密。每个省都成立统一领导、指挥地方武装的军区，直辖数个独立师和军分区；每个县均有独立团，几个县合组为一个军分区。军区的独立师和军分区的独立团作为半正规性的地方武装，具有一定战斗力。各县除独立团外，还有以保卫地方为主要任务的非正规性军事组织，边区县有警卫连、营和游击队，腹地县有赤卫军。此外，县、区、乡尚有半军事性的少年先锋队和童子团。红军—地方独立师—赤卫队—少先队，苏区武装形式形成一个有机链条，既可以帮助保证武装力量的人员供给，又可以在持续的战争状态下，以强有力的军事形式实现对社会的凝聚和控制。更重要的，如论者所言，这一制度：

\begin{quote}

将一村一乡的自卫队，联结为一张遍及苏区的大网，将历来难以守卫自身利益的个体农民，联结为一股股有组织、有指挥的武装力量，从而在地方武装中赋予了组织民众、武装民众的社会意义和战略意义。完全反映了这种思想的苏区地方武装建设，在苏区社会创建和中国革命发展中，首开人民武装制度和全民皆兵制度之先河。\footnote{何友良：《苏区制度、社会与民众研究》，未刊稿，第121页。}
\end{quote}


苏区地方武装的成员几乎包括除老弱婴幼及阶级异己人员以外的所有群众，涵盖面居苏区社会各种组织之首。以性别分有男、女赤卫队；以年龄分，24岁至50岁者加入赤卫军、赤卫队，16岁至23岁者加入少年先锋队，8岁至15岁者加入童子团。从地方武装所占人口比例看，如兴国县长冈乡，赤卫队年龄段男子全乡共66人，除乡政府主席、文书及重病残疾者共20人外，共余46人全都加入了赤卫队；全乡女子246人，除病残的26人外，全部加入女子赤卫队；全乡101名少年先锋队年龄段的男女，也仅有病残的15人未入队。据江西省1932年的一个统计，“全省赤卫军的组织至少有二十五万，连半军事性的少队组织至少有五十万”。\footnote{《江西苏区中共省委工作总结报告（1932年5月）》，《中央革命根据地史料选编》（上），第442页。}大多数赤卫军和少先队建立了经常性的军事操练和政治训练制度。赣西报告：“群众的军事训练有相当好，特别是少队童团要好些。有一次少共举行西路少共童团总会操（期在五卅），永新、安福西区团江区等都有少队童团比赛起来，确实操得整齐特别精神好，动作也操得十余种，尤其是儿童团操得更好。据说赤军较少队还要逊一着，在政治测验童团也要好”。\footnote{《张启龙关于赣西工作情形给中央局的报告（1931年7月12日）》，《江西革命历史文件汇集（1931年）》，第102页。}兴国、赣县、胜利、上杭等县的许多区，涌现出赤卫军模范营，能够直接参加作战并能调到别的地方配合红军或独立师团作战。1934年4月28日，中革军委总动员武装部公布群众武装统计结果是：赤卫军人数共266345人，少年先锋队共157149人，模范赤卫军共55695人，模范少年先锋队共36273人。\footnote{李光：《中国新军队》，1936年编印，第227～229页。李光即滕代远。}


对苏区群众性武装的作用，毛泽东曾代表苏维埃中央予以高度肯定：“他们之加入红军，他们之保卫地方，他们之袭敌扰敌，在历次粉碎‘围剿’的战斗中，显示了他们极其伟大的成绩，致使敌人惊为奇迹，成为敌人侵入苏区的绝大困难。”\footnote{毛泽东：《对第二次全国苏维埃代表大会工作报告》，《苏维埃中国》，第257页。}作为对手方的蒋介石对苏区的组织状况也颇为艳羡，他在对属下的讲演中称赞苏区：

\begin{quote}

政治组织和民众的组织，都很严密，尤其是民众的组织，我们最不及他。匪区的民众，他们都尽量的组织并武装起来，成为各种的别动队，如赤卫队、慰劳队、童子团、少年先锋队等都是，他们因为组织比较来得普遍而严密，所以民众生活能军事化，一切行动能够秘密灵敏。尤其是他的侦探来得很确实、很快速，通讯网递步哨号炮以及其它种种侦探的标志和手段，统统都能奏效。其实他的侦探和通讯联络的方法，譬如放几个号炮，摇几下红旗子，都是极简单的，但是因为得力于组织的严密，却能格外迅速确实而发生很大的效用，使匪区民众和他们伪政府匪军，都能协同一致动作。\footnote{蒋介石：《剿匪基本工作之研究》，《先总统蒋公思想言论总集》第11卷，第37～38页。}
\end{quote}


蒋介石上文中说到的通讯联络方法，康泽曾有详尽的报告：“伪军于驻地各村庄设立通讯站，预于两站之间择定置讯处所将信筒埋于砖石下丛草中，利用老幼妇女假装出外工作辗转传递极为敏捷。且不易被人察觉。伪军侦探多使用老人幼童妇女，活动时各备划有痕迹之铜币信件为的号，其携带信件常缝入衣裤或折于袖口上。”\footnote{康泽电蒋中正伪军于驻地各村庄设立通讯站两站之间择定置讯处所及伪军侦探多用老人幼童妇女传递信件，蒋中正文物档案002090300085199，台北，“国史馆”藏。}在当时通讯条件不发达的状况下，这样的讯息传递方式可谓巧妙，而民众参与其中，如果没有精密的组织及发自内心的支持，都是难以想象的。


可以证明苏区组织严密的具体事例不胜枚举，当时报刊记载苏区查验路条的情况：

\begin{quote}

有一天早晨，我在离于都城不远的麻油坑地方经过。那时天色才黎明不久，路上行人还很稀少……突然，尖锐的问话传来：“同志，有路条么？”我抬头一看，原来是一位青年女同志。我回答她：“我没有路条。”“你没有路条，不准通过；我们要把你扣留起来，送到乡苏去。”她一面说，一面前来检查我。……最后我把介绍信给她看，因为她不认得字，连忙叫男同志来看，一个同志看了还不算，再给第二个同志看，证实了才准我通行。\footnote{郭庆福：《模范的女性》，《青年实话》第2卷第7期，1933年3月12日。}
\end{quote}


总的来看，这一套控制体系设计非常严密，国民党方面相关调查谈道：“边匪地方，以五里一哨，中心区域，则十里一哨。对于盘查，十分认真，如无匪区县政府之路条，则行扣留。凡是群众行动，必须有高级军事机关护照，方能通行，并于每月终，由县军事部派人检查各区步哨一次，各区军事部，每十天须派员检查一次。”\footnote{《赣西匪区之近状》，《军政旬刊》第2期，1933年10月30日。}“民众之往来均须有路票始能出入，管束极为严密。若开民众大会即由代表率领前去。”\footnote{《石凌生口供》，《西路军公报》第9期，1934年3月15日。石被捕前为中共龙武县委书记。}


当然，习久成疲，苏区步哨虚应故事的情况也不是没有：

\begin{quote}

我于3月6号，由博生到太雷做突击工作，经过一百余里的路途，不曾见有一个步哨，最后我由太雷县委于10号出发经日车，到龙岗区的一乡，才见一个步哨。这个同志在店内问：“同志有介绍信么？”我即答复：“没有。”他说：“没有过不得。”我不去管他，他就坐在凳上，再也不问我了。我走过以后，我到回去问他：“同志你在这里做什么？”他答：“放哨呵。”\footnote{《不让一个反革命分子漏网》，《青年实话》第3卷第15期，1934年3月。}
\end{quote}


有意思的是，当时国民党方面材料对中共组织严密度的估计甚至要高于中共本身，这一方面是对对手的重视，另一方面也是对自身缺陷的明确认知；而对中共而言，承认并努力发现自身的不足，正是其继续改进、提高自身的重要动力。


作为权力结构深入乡村、社会动员能力空前强大的政权，苏维埃政权有能力在社会救济方面发挥比传统社会更多的作用。苏区专门成立互济会，负责展开互助救济，乡设委员会，村有负责工作的主任，下再分小组，互济会成员包括大部分苏区民众，每人每月交会费一个铜板。毛泽东在兴国长冈乡的调查记载了第五次反“围剿”期间社会救济状况：

\begin{quote}

募捐救济难民，援助反帝。今年有过二次。一次是七十多个信丰难民到兴国城（榔木乡时），共捐了二十多串。一次是援助东北义勇军（也是榔木乡时，那时人口二千九百，会员约八百），捐了四十多串。捐数五个铜片起，一百的、二百的、一串的都有。一百的多数，约占会员百分之六十。五个铜片的，一串的，各只几人。

（三）乡里火烧了房子的，失业工人生病无药的，募捐救济。今春一家失火，烧了一间半屋，捐了六串多钱给他。

（四）救济饥荒。今夏榔木乡有三四个人饿饭（过去乞丐，现还很穷），请求区互济会发钱发米救济，每人每次多的三升，少的一升，今夏发了三四次。\footnote{毛泽东：《长冈乡调查》，《毛泽东农村调查文集》，第322页。}
\end{quote}


这些救济活动虽然琐碎，但对患难中的乡民意义却非同一般，有助于增进群众对苏维埃政权的信任。不过，作为群众性组织，互济会的活动也有尴尬之处，一方面，根据毛泽东的说法：“许多地方的苏维埃不注意社会救济工作，许多地方的互济会只知收月费不知救济群众困难。”\footnote{毛泽东：《长冈乡调查》，《毛泽东农村调查文集》，第323页。}另一方面，当互济会真正发挥作用时，由于其和民众利益息息相关，又发生“苏维埃召集群众会议而群众不到会，互济会召集会议则多数群众到会”的情形，导致地方苏维埃政权“简单的将互济支部取消”。无论哪种情况，最终的结果都导致互济会空心化，使其“组织上失掉一个下层系统，一切工作都陷于停顿状态，同时上层机关庞大，找不到实际工作去干，成为机关运动的形式”。\footnote{《互济会问题决议案》，《湘鄂赣革命根据地财政经济史料摘编》（下），第938页。}以致有些地区互济会“对于救济工作根本没有做到，因此，群众对互济会的经费由怀疑而发生不满，甚至说互济会是捐税机关”。\footnote{《中国革命互济会湘鄂赣省总会筹备处通知》，《湘鄂赣革命根据地财政经济史料摘编》（下），第941页。}更有一些地区互济会出现侵吞现象：

\begin{quote}

募捐手续，不能用宣传鼓动方式，仅采取摊派式的要群众照数缴钱的办法，而收入支出之数，又不能按月公开宣布，其中浪费、滥用、舞弊都在所难免。兹仅据修水一县的互济总会报告，自去年七月至今年十月共收入银洋八千六百四十多元，其中用在救济事业方面的只有一千七百多元，用在伙食、零用、办公、长驻救护队等方面的，达到一千九百多元，此外是县委挪用一千多元，县苏挪用三千多元，尤其是私人挪用一百五十多元，竟有半数（七十多元）无从收回的。\footnote{《互济会问题决议案》，《湘鄂赣革命根据地财政经济史料摘编》（下），第938页。}
\end{quote}


可以看出，制度设计和实际操作间还有一段距离。


对苏区内部一些游离力量的改造、重组是完善社会组织的重要任务。闽赣两省游民众多，据1930年苏区有关文件的分析，游民无产者包括多种从事“不正当职业”者，具体有：盗贼、娼妓、兵痞、戏子、差人、赌棍、乞丐、人贩子、巫婆等。对这些人，中共采取改造为主方针，相当一批游民在苏维埃革命中分得了土地，生活习俗得到改造，成为自食其力的普通劳动者，有的还进入苏维埃政权。如兴国永丰区共有游民90多人，他们在革命中得到很多利益，因此，“一般都是欢迎革命，不但没有一个反革命的，并且有十个参加区乡政府的指导工作，一个当了游击队的指挥员”。\footnote{毛泽东：《兴国调查》，《毛泽东农村调查文集》，第233页。}


不过，并不是所有的流民都顺利地接受了改造，苏区时期一直没有得到解决的会匪问题就颇让人伤神。闽赣两省会党、土匪活动频繁，中央苏区所在的赣南及闽西会党组织和土匪尤为活跃。据中共方面估计，闽西各地游民“占当地人口的百分之二十五”，\footnote{《中共闽西第一次代表大会之政治决议案》，《中央革命根据地史料选编》（中），第111页。}比例相当惊人。这些游民当中，又有相当一部分和土匪难以分开。赣南、闽西山高林深，宗族组织强盛，这些都便利于土匪成形，民国以来中央权威混乱、地方控制软弱，又为土匪大规模孳生提供了条件，所以，民国年间，这里土匪活动一直旺盛。地方政府常常抱怨对土匪“地广山多，林深将密，包围搜捕，两穷于术。”\footnote{《李厚基请奖镇压德化等县地方变乱军警文职人员咨（1916年9月28日）》，《民国时期泉州档案资料选编》，中国第二历史档案馆、泉州市地方志编委会编印，1995，第124页。}


按照中共在苏维埃革命初期的理解，会党属于游民无产者的一部分，因此在早期的暴动中，发动会党参加是暴动的重要一环，中共文件指示：“应特别注意两点：一、我们要会匪与农民建立兄弟般的关系，不要存利用的心理，不是希望他们来帮助，而是要他们与农民一样为暴动的主体；二、我们联络会匪要拉住群众，不要仅拉他们的领袖。”\footnote{《中共江西省委秋收暴动计划》，《江西革命历史文件汇集（1927～1928）》，第21页。}国民党方面观察：“江西为洪江会策源地，八十一县，几无一县不有洪会机关……赤匪知之，遂利用洪匪为内奸，标语欢迎洪家兄弟，互相勾结，与匪打成一片。凡赤卫军梭标队皆以洪江会匪充之。盖今日之所谓农匪，无一不是洪江会匪也。”\footnote{《刘纯仁条陈》，《陆海空军总司令部行营党政委员会会报》第5～8期合刊，1931年8月31日。}蒋介石也曾要求各地对“会匪曾扰害人民或勾通赤匪查有确实证据者，并仰按名拿案，尽法查办，以资儆惕。其有确被裹胁或盲从附和者，得勒令缴出票布，从宽免究，予以自新。务期境内永无洪会之非法团体与洪会之匪徒”。\footnote{《训令各县会准行营参谋处函送刘纯仁条陈请查照核办》，《陆海空军总司令部行营党政委员会会报》第5～8期合刊，1931年8月31日。}


苏维埃政权成立后，会匪作为游离于正常秩序之外的武装力量，和苏维埃政权必然发生冲突。中共针对会匪制定了很多政策，并采取措施要求“坚决执行阶级路线与充分的群众路线，号召被欺骗压迫误入刀团匪的份子回来分田，而且要在群众中详细解释”，\footnote{式平：《查田运动与肃清刀团匪》，《红色闽赣》第4期，1934年2月13日。}但面对不断到来的“围剿”，不可能有更多时间、精力彻底解决会匪问题，会匪大有卷土重来之势。而且随着阶级对立的加深，一部分地主出身者铤而走险也加入到会匪行列。1934年前后，由于苏维埃政治经济政策的失误，更客观上壮大了会匪队伍。当时在苏区的一些山高谷深地区，会匪活动尚很猖獗。福建长汀有“李七孜、曹丰溪、马贤康等十余股零星残余团匪土匪……四百多枝枪”。\footnote{《福建省苏第二十四次主席团会议决议案》，《中央革命根据地史料选编》（下），第213页。}广昌县苏国民经济部部长郑会文“去年讨老婆受了大刀会匪首领的贺礼二元，后来他的老婆被另一部分大刀会捉去”，\footnote{《江西省苏维埃国民经济部命令，第4号（1934年2月22日）》，陈诚档案缩微胶卷008·638/0067/0152，中国社会科学院近代史研究所藏。}刀匪的胆大妄为非同小可。1933年在红军总政治部教导队学习的肖光多年后仍清楚记得自己当年的亲身经历：

\begin{quote}

总部和前委在建宁时都驻在一起。记得建宁地方大刀会较多，红军站岗怕他们来摸哨，有一个晚上，我在城墙边站岗，看到一团黑呼呼的前来，叫了口令没回答，认为是大刀会来摸哨，就开枪打，一看是打死了一条牛。教导队只好给群众赔了钱。\footnote{《访中央档案馆肖光同志纪实（1978年7月30日）》，《建宁党史资料》第2辑，中共建宁县委党史办公室编印，1987，第2页。}
\end{quote}


中央苏区地方政权及武装和土匪的较量一直延续到中央红军的撤离。1929～1934年10月，宁都县的赤卫队和游击队等与地主武装、靖卫团、大刀队等有过200多次战斗。\footnote{《宁都县志》，江西人民出版社，1986，第399页。}土匪武装不仅仅消极躲避以求生存，时有主动挑衅，福建泰宁“大刀会经常下山袭击我们的机关或者在中途打我们的埋伏……泰宁每个区都被袭击过，区委大部分牺牲了，朱口区有一次被包围，只剩下一个同志，其它都牺牲了”。\footnote{《苏州军分区政委谭成章同志口述记录稿》，《泰宁文史资料》第1～3辑合订本，第25页。}赤白边境地带红军后方运输也是其骚扰对象：“东方军兵站线，由建宁经大田市、新桥、下南、上关、飞鸢或纸马到前方，沿途刀匪不安。”\footnote{《朱德、周恩来致项英电（1933年10月6日）》，《中央苏区第五次反“围剿”》（上）（《江西党史资料》第21辑），中共江西省委党史资料征集委员会、中共江西省委党史研究室编印，1992，第83页。}甚至对红军大股部队也展开袭击，1934年7月，北上抗日先遣队由江西出发进入福建山区后，就常常遭到大刀会袭击：“这些家伙，用朱砂涂成大花脸，红花绿袍，身上贴着鬼符，张牙舞爪，猛看一眼倒挺吓人。他们凭借着地形熟悉，一会儿漫山遍地都是，敲锣打鼓偷袭我们，一会儿几乎全不见了。”\footnote{谭志刚：《记中国工农红军北上抗日先遣队》，《江西文史资料选辑》第3辑，政协江西文史资料委员会编印，1980，第51页。}


刀匪的活跃，其实还提示着中央苏区的另一种现实，即在崇山峻岭、村庄散落的赣南、闽西地区，即便像中共组织力如此强大的政治力量，短时期内也有力有未逮之处。当年红军的日记记载，1934年红九军团在广昌的头陂时，到离头陂20里的一个村挑土豪的存谷，头陂周边百里此时都是红色区域，该地理应在苏维埃政权控制下，但日记所载挑谷经过是：“侦察员先进村侦察，武装部队随后跟进，严密警戒。那里的青壮年大部分被国民党抓去当兵了，村里多系老弱妇孺替他们管粮仓，地主豪绅都搬进了大城市，所以运谷很顺利。”\footnote{赵镕：《长征日记》，1934年4月21日，山西人民出版社，1990，第64页。}苏维埃革命数年后，在已属中央苏区核心地区的头陂，还有大量的土豪存谷，而从红军进入该村的状态看，该村和革命几乎尚处于绝缘状态。这样的事例的存在，再次显示着一旦进入到历史的具体情境中，事态总会比想象的更加复杂。虽然这并不足以动摇对苏区组织基本状况的判断，但却提醒我们，许多原则性的结论后面，或许都还存留着更多可能出现的丰富细节，历史的弹性常常超乎人们的想象。

\section{红军：堡垒的坚强核心}


作为武装割据的军事根据地，苏区具有十分明显的军事特征。尤其在早期，由于红军的弱小，对固定区域实行长期占据的力量不足，苏区区域经常随着红军的流动而变化。因此，在苏区建设中，红军常常担负起苏维埃政权的作用，成为中共和民众之间联结的纽带。红军严明的纪律、良好的政治素质使其与民众间建立了十分亲密的关系。这种关系在国民党方面的材料中处处可见，兹引几例，以见一斑。1931年第一次“围剿”失败后，国民党军所编战报称：“匪区民众，久受赤化，所得我方消息，即行转告匪军。”\footnote{《关于第一次赣南围剿之经过情形》，《中华民国史档案资料汇编》第5辑第1编军事（3），第50页。}第十八师报告则写得更详尽：

\begin{quote}

东固暨其以东地区，尽属山地，蜿蜒绵直，道路崎岖，所有民众，多经匪化，且深受麻醉，盖匪即是民，民即是匪。对于我军进剿，不仅消极的认为恶意，且极端仇视，力图抗拒。如是，对于我军作战上发生下列之困难：（一）我军师行所至，农匪坚壁清野、悉数潜匿山中；（二）潜伏山中之匪徒，对于我军状态窥探无遗。如是，我军企图完全暴露；（三）我军不仅不能派遣一侦察，即欲寻一百姓问路，亦不可得，以故我方对于匪情全不明确。即友军之联系亦不容易；（四）山地道路崎岖，行军已感困难，而匪徒对前进之道路亦无不大加破坏，我之前进，几使我无路可走，盖一则可予我之极大疲劳，一则无形中可迟滞我军。\footnote{《第十八师失败之检讨》，《东固·赣西南革命根据地史料选编》（下），中央文献出版社，2007，第769页。}
\end{quote}


参加第三次“围剿”的蔡廷锴，日后回忆其进入苏区后的情况时说道：“地方群众在共党势力范围下，或逃亡，或随红军行动，欲雇挑夫固不可能，即寻向导带路亦无一人，至于侦探更一无所得，变成盲目。”\footnote{《蔡廷锴自传》（上），黑龙江人民出版社，1982，第245页。}


红军组织之严密、对国民党军研究之精细有时简直让人惊叹，朱德后来在和美国记者史沫特莱谈话时，谈到红军的情报工作：

\begin{quote}

红军的情报工作当时已经组织得很好，不但满布苏区，而且深入国民党区域。红军开设了特别训练班来训练情报人员，内中有不少是妇女和儿童，另有一些小贩和串村子的手艺人，他们的工作便于在国民党区域内广泛活动，可以到有钱人家或穷人家做工或卖东西，也可以混入敌军营盘内活动。

……

有一组人专门研究敌人的电码、公报、出版物，并且和俘虏谈话。另一组人负责从新占领区收集情报。再有一组人专搞历史工作——仔细研究敌方每一个军，把军官和士兵的背景都调查清楚：这个军是从哪一省来的，军里有什么变化，过去的历史和组织，它的战斗力，等等。根据这些研究，朱将军解释道，“我们最后可以决定应付某一个军的最好办法。”\footnote{〔美〕史沫特莱：《伟大的道路》，梅念译，三联书店，1979，第346页。}
\end{quote}


如此用心深密的研究、组织，的确不能不让其对手方不寒而栗。


红军的成功建设当然和中共确立的一整套政治、思想和组织原则密切相关。从三湾改编到古田会议决议，中共提出和坚持党对军队的绝对领导，通过将支部建在连上等一系列政治组织措施保证党领导枪原则的贯彻实施，严防军队领导权落入个人或小集团手中。同时，规定全心全意为人民服务为建军的唯一宗旨，把军队的发展同人民群众的根本利益相联系，古田会议决议明确指出：

\begin{quote}

中国的红军是一个执行革命的政治任务的武装集团。特别是现在，红军决不是单纯地打仗的，它除了打仗消灭敌人军事力量之外，还要负担宣传群众、组织群众、武装群众、帮助群众建立革命政权以至于建立共产党的组织等项重大的任务。红军的打仗，不是单纯地为了打仗而打仗，而是为了宣传群众、组织群众、武装群众，并帮助群众建设革命政权才去打仗的，离了对群众的宣传、组织、武装和建设革命政权等项目标，就是失去了打仗的意义，也就是失去了红军存在的意义。\footnote{《中国共产党红军第四军第九次代表大会决议案》，《中共中央文件选集》第5册，第801页。}
\end{quote}


红军对民众的友爱表现，自然得到民众的回报。红军受到民众欢迎，民众表现出罕见的参军的积极性。送儿、送夫参加红军的事迹在苏区常常可见。福建上杭官庄区余坊乡林东姑“不但宣传自己的儿子去当红军，而且宣传别人叫儿子去当红军，对人家（指被宣传者）说：‘你的儿子能去当红军，群众同（给——引者注）我作冲锋劳动工、同我砍的柴都给你。’”。\footnote{李中：《“复光运动周”中上杭劳苦青年的光荣》，《青年实话》第2卷第19号，1933年6月18日。}旧县区新坊乡的李永春，“父子两人一起报名当红军”。\footnote{李中：《“复光运动周”中上杭劳苦青年的光荣》，《青年实话》第2卷第19号，1933年6月18日。}长汀的吴秀英为动员自己的未婚夫当红军，主动与其结婚，然后动员其加入少共国际师。\footnote{唐铁海：《中央老根据地印象记》，劳动出版社，1952，第68页。}妇女为支持自己的亲人参军，甚至不惜以离婚相威胁，汀东长宁区彭坊乡的江银子，宣传丈夫当红军，“她的丈夫不去，她就到乡苏要求同她丈夫离婚，后来她的丈夫就自动报名当红军了”。\footnote{《光荣的红匾上》，《青年实话》第2卷第19号，1933年6月18日。}瑞金也有报告：“桃黄区有一个妇女要他老公去当红军，他不，就向他离婚。”\footnote{《中共瑞金县委九月份工作综合报告》，《江西革命历史文件汇集（1932年）》（二），第142页。}


虽然随着苏区的发展，国共之间“围剿”与反“围剿”战争深入持久进行，战争的残酷性凸显，加上其他一些原因，民众参军的热情在逐渐下降，但处处仍可看出红军受到苏区人民的真心支持：“当红军与敌人持久战的时候，广大群众送饭到火线上去给红军吃。在送饭、送菜的时候，路上是成千成万的饭担子，群众送的饭，红军主力还吃不完，并且群众把自己的油盐、瓜菜不吃，都弄去送给红军吃。”\footnote{《关于湘鄂西具体情形的报告（1932年12月）》，《湘鄂西苏区革命历史文件汇集》第1册，第312页。}1932年初，福建“永定溪南区群众听说五军团入闽，挑许多饭菜在龙岗等了一天”。\footnote{《团闽粤赣省委三个月工作报告》，《闽粤赣革命历史文件汇集（1932～1933）》，第160页。}红军和民众之所以能够保持如此亲密的鱼水关系，关键在于其苏维埃革命的政治目标和严格的政治、军事纪律；同时，红军的组成更牢固了这一关系。据1934年4月红军总政治部统计，中央苏区红军中有77%的人员都来自中央苏区本身，其中工人、农民成分占98%。\footnote{李光：《中国新军队》，第277～278页。}这种苏区民众子弟兵的特色，使民众和红军具有了天然的紧密联系，既增进了民众对红军的信赖与爱护，也有利于军队纪律的执行，这是外来军队所难以比拟的。同时，红军始终注意和民众利益相通，打土豪时往往将收获的一部分分给群众，在大多数情况下，“群众都很热烈的要来参加分东西”。\footnote{《北冀关于漳州红三团行动的报告（1933年8月25日）》，《福建革命历史文件汇集》甲11册，第89页。}


红军具有良好的组织和宣传系统，十分注意对官兵进行政治和文化教育，有一系列提高官兵政治和文化素质的措施。陈毅给中共中央汇报中，谈到了红四军的政治训练：

\begin{quote}

（一）讲演，由官长召集全体讲话，或作政治报告，或作生活批评，或作工农运动概况报告等。

（二）讲课，在军队有三日的休息，作每日必有一小时政治课，由党代表担任去讲，这个讲演比较有计划的，或定于一月讲演许多题目，这些题目是可以连贯的使士兵得到一些有系统的政治智识。

（三）早晚点名讲话与呼口〈号〉，则例每日士兵生活批评或对于明日行动之煽动宣传等。

（四）在一次游〈击〉工作，一次战斗，一次行动，经过以后的批评，要详细向士兵讲出来。

（五）军队里举行识〈字〉行动，简易的办法就是要士兵认红军的标语，认得一个标语即将此标语包含的意义策略告诉他。\footnote{《陈毅关于朱毛军的历史及其状况的报告（1929年9月1日）》，《中共中央文件选集》第5册，第760页。}
\end{quote}


这些措施着眼于提高官兵政治思想水平，顾虑周到、设计细致，无怪乎周恩来主持的《军事通讯》在刊载陈毅的报告时加编者按说：这里面有很多宝贵的经验值得我们每一个同志注意，如他们的编制，他们的战术，他们的筹款给养的方法，他们与群众的关系，他们对内的军事和政治训练，他们处置军中给供开支的原则（官兵夫经济平等，开支能公开）……都是在中国“别开生面”，在过去所没有看过听过的。


红军利用各种条件加强政治教育。红三军团“十一团的墙报平均三天能够出版一次，投稿人数最多的有八十多人，最少的也有五十多人投稿”。\footnote{三军团政治部《武库》第8期，1932年1月12日。}红军每个连队都设置列宁室，作为政治教育和文化学习的场所，曾经和红军一同行军的外国传教士写道：“每到一地，不管时间长短，‘列宁室’是必建的。所谓‘列宁室’，实际上就是红军读书学习的一个地方，有时利用房子，有时就自己动手临时建，8根竹竿或树杆做柱，绿色的树枝和竹枝编在一起做墙，屋顶铺上稻草就算天花板。这个地方就成了他们看书学习或集体活动的地方了。”\footnote{《外国人笔下的中国红军》，陕西人民出版社，1996，第261页。}通过这些持续不断的宣传教育，使红军成为和中国历史上其他军队有很大不同的一支武装。天津《益世报》报道中总结了红军的三个优点：

\begin{quote}

一、官长兵士均有政治意识，故能严明军纪，指挥如意，尝以镜窥其阵线，见彼等虽在战壕中，亦时有开会演说之模样，可知军事政治化之运动，虽在战时，亦不放松；二、注意宣传，该匪军近水楼台，日受麻醉宣传，固无论矣，唯彼且将各种标语或用纸印，随风飘送，或书木片，顺水漂出，以图赤化良善军民，影响不少；三、应战沉着，彼匪若取守势，无论国军如何攻击，均不轻易还击，以致难窥虚实。\footnote{《军事中心的临川》，1933年8月4日《益世报》。}
\end{quote}


上述三点准确道出了红军组织纪律的严明及对宣传教育的重视。这些教育，使红军官兵了解他们的使命、目标及与群众的血肉联系，促进其自觉遵守纪律，保持与群众的良好关系。同时，红军建立严密的监督体制，成立纪律检查委员会，保证纪律的有效贯彻。当年红军的日记中记有纪律检查的实际状况：

\begin{quote}

各单位纪律检查委员会对部队离开广昌的情况进行了联合检查，到新安宿营后进行了评比：军团政治部第一，司令部第二，供给部第三，惟军团部的警卫连、侦察连做得不好，大多住户的水缸没挑满，房前屋后的垃圾也没清除。供给部各科的住地本来打扫得很干净，给老百姓的水缸也挑满了水，但发现运输队有人买老百姓的鸡蛋价钱不公平，有一个队把老百姓扁担损坏了没赔偿，有的铺草没捆好送到原地，门板送错了也不问。这些问题反映了民夫们还不习惯于严格遵守三大纪律八项注意。\footnote{赵镕：《长征日记（1934年4月13日）》，第60页。}
\end{quote}


严格的纪律要求不仅针对红军，甚至及于随军的民夫，其严厉程度不由得让人惊叹。对此，蒋介石不得不自叹不如：“讲到军纪方面，土匪因为监督的方法很严，无论官兵，纪律还是很好，所以在战场上能勇敢作战，而对于匪区一般民众，还是不十分骚扰。我们的情形老实说起来，是不如他们！”\footnote{蒋介石：《剿匪技能之研究》，《先总统蒋公全集》第1册，第674页。}


良好的宣传教育及思想政治工作不仅造就了红军严密的组织、良好的群众关系，更重要的，对于一支武装力量而言，还使其获得了坚强的战斗力。在与拥有资源、给养、人员、武器等多方面的绝对优势的国民党军的较量中，红军屡屡以弱胜强，就是这一点的最好证明。


当然，正像任何事物都不可能完美无缺那样，红军的发展过程也不是一马平川，其间难免有曲折、坎坷，如中共中央领导人进入苏区后，红军发展就出现一些波折。


应该说，中共中央领导人到达苏区，对红军的影响也有正反两方面。积极方面，毛泽东曾总结出五点：

\begin{quote}

（一）成分提高了，实现了工农劳苦群众才有手执武器的光荣权利，而坚决驱逐那些混进来的阶级异己分子。（二）工人干部增加了，政治委员制度普遍建立了，红军掌握在可靠的指挥者手中。（三）政治教育进步了，坚定了红色战士为苏维埃斗争到底的决心。提高了阶级自觉的纪律，密切了红军与广大民众之间的联系。（四）军事技术提高了，现在的红军虽然还缺乏最新式武器的采用及其使用方法的练习，然而一般的军事技术，是比过去时期大大进步了。（五）编制改变了，使红军在组织上增加了力量。\footnote{毛泽东：《中华苏维埃共和国中央执行委员会与人民委员会对第二次全国苏维埃代表大会的报告》，《苏维埃中国》，第257页。}
\end{quote}


毛泽东这一总结是1934年代表临时中央政府所作，虽然不一定完全符合他的本意，但大体还是反映了当时的实况，不能认为完全是言不由衷之词；而就消极方面言，1940年代中期撰写的《在七大的建军报告》初稿对此有所论列：“在军民关系上取消了红军作群众工作的优良传统，只赋予红军以单纯的作战任务。在动员参战工作上，不懂得培养和爱惜苏区的人力物力财力以供长期战争的适当调用，而是用的竭泽而渔的办法，在两年之内便把苏区的力量弄到枯竭的地步，迫使红军不得不退出苏区去另寻生路。”\footnote{《在七大的建军报告》（初稿），1945年3月1日。}这一论述触及问题的某些症结，是中共善于自我批评、力求上进的体现。不过，如果从更广泛的角度看，苏区力量枯竭的责任不能仅仅算到中共自己头上，苏区本身资源的困乏、国民党方面的封锁与逼迫事实上更是肇祸元戎，中共中央指导上的问题只是在内外环境不利的背景下才被放大，进而一定程度上动摇到中央苏区党、政权、红军的坚强基础。从我们将要面对的第五次“围剿”与反“围剿”进程看，当时苏区某种程度上的“师老兵疲”，既不仅仅是中共本身的问题，也非突然凭空而降，实为苏区在内外挤迫下遭遇困境的集中体现。

